%% Sakarya Universitesi Fen Bilimleri Enstitusu - Lisansustu Tez Sablonu (LaTeX)
%% Derleme:  latexmk -lualatex tez.tex      (veya: lualatex tez ; lualatex tez ; lualatex tez)
%% Secenekler: doktora | yukseklisans , notlar | notsuz (kirmizi sablon aciklamalari),
%%             wordline (Word satir modeli, varsayilan) | texline (klasik TeX satir kirma),
%%             wordbirebir (Word sablonunun altbilgi farkliliklarini da kopyalar: 1. bolumdeki yuksek
%%             sayfa numarasi, yatay sayfalarda koyu olmayan numara; kaldirilirsa tum numaralar ayni hizada)
\documentclass[doktora,notlar,wordbirebir]{sauthesis}

%----------------------------------------------------------------- tez bilgileri
\tezbasligi{TEZ BAŞLIĞI BURAYA YAZILIR\\GEREKLİ İSE İKİNCİ VE ÜÇÜNCÜ SATIR KULLANILABİLİR\\ ÜÇ SATIRA SIĞDIRINIZ}
\tezbasligitr{TÜRKÇE TEZ BAŞLIĞI BURAYA YAZILIR}
\tezbasligien{THESIS TITLE IN ENGLISH IS WRITTEN HERE}
\tezbasligibeyan{TEZ BAŞLIĞI YAZILACAK}
\teztipi{YÜKSEK LİSANS/DOKTORA TEZİ}   % "DOKTORA TEZİ" veya "YÜKSEK LİSANS TEZİ"
\yazar{Öğrenci Adı SOYADI}
\yazarbeyan{Öğrencinin Adı Soyadı}
%\anabilimdali{Bilgisayar Mühendisliği}   % bos birakilirsa gri "Bir öğe seçin." basilir
%\bilimdali{...}
\danisman{Unvan Ad SOYAD}
\ortakdanisman{Unvan Ad SOYAD}
\kapaktarihi{AY YIL}
\savunmatarihi{16.06.2022}
\tesekkurimza{Ad Soyad}

\begin{document}
%----------------------------------------------------------------- on kisim
\diskapak          % dis kapak + arkasi (numarasiz)
\ickapak           % ic kapak (i) + arkasi (ii) (numara yazilmaz)
\onaysayfasi       % iii
\clearpage\null\thispagestyle{plain}\clearpage  % iv (bos)
\beyanname         % v
\clearpage\null\thispagestyle{plain}\clearpage  % vi (bos)
\ithaf{Eşime ve çocuklarıma (Örnektir, farklı bir ifade yazılabilir)}   % vii
\begin{tesekkur}
Teşekkür bölümünün içerisindeki metinler 1 satır aralıklı, önce ve sonra 6 punto (nk) aralıklı olarak yazılır. Tezin ilk sayfası niteliğinde yazılan teşekkür bir \textcolor[HTML]{FF0000}{}sayfayı geçemez.

Tez danışmanına (ve varsa ortak danışmana), tezi destekleyen kurum ve kuruluşlara ve yardımcı olan kişilere bu kısımda teşekkür edilir. Teşekkür metninin altında sağa dayalı olarak tez sahibinin adı ve soyadı yazılır.
\end{tesekkur}

\tableofcontents
\begin{kisaltmalar}
\kisaltma{AIC}{Akaike Information Criteria}
\kisaltma{BHA}{Bütillendirilmiş hidroksi anisol}
\kisaltma{BHT}{Bütillendirilmiş hidroksi toluen}
\kisaltma{DPPH}{1,1-Difenil-2-pikrilhidrazil}
\kisaltma{FRAP}{Ferric Reducing Antioxidant Power}
\kisaltma{GAE}{Gallik asit eşdeğeri}
\kisaltma{LDL}{Düşük yoğunluklu lipoprotein}
\kisaltma{ROT}{Reaktif oksijen türleri}
\kisaltma{TBARS}{Tiyobarbutirik asit reaktif maddeleri}
\kisaltma{TPTZ}{2,4,6-Tripiridil-s-triazin}
\kisaltma{TSA}{Tyriptic Soy Agar}
\kisaltma{TSB}{Tyriptic Soy Broth}
\kisaltma{ÜÇE}{Üzüm çekirdeği ekstraktı}
\end{kisaltmalar}

\begin{simgeler}
\simge{C}{Dokunun kapasitansı [Birim]}
\simge{H}{Isı miktarı [Birim]}
\simge{q}{Faz yükü [Birim]}
\simge{t}{Zaman [Birim]}
\simge{u,v}{Yer değiştirme vektörü bileşenleri}
\simge{w}{Açısal hız [Birim]}
\simge{XC}{Kapasitif reaktans}
\simge{XL}{Endüktif reaktans}
\simge{α}{Asal gerilme doğrultusundan sapma açısı}
\simge{ρ}{Yoğunluk [Birim]}
\simge{D}{Çap [Birim]}
\simge{ρ}{Yoğunluk}
\simge{m}{Kütle [Birim]}
\simge{μ}{Dinamik Viskozite}
\simge{E}{Elastisite Modülü}
\simge{ν}{Poisson’s Oranı}
\simge{b}{Bulk Modülü}
\simge{\ensuremath{\symbfup{σ}}}{Couchy Gerilme Tensörü}
\simge{ε}{Şekil Değiştirme Tensörü}
\end{simgeler}

\listoftables
\listoffigures
\begin{ozet}
\notrenk{FF0000}{Özet hazırlanırken 1 satır boşluk bırakılır. Türkçe tezlerde, Türkçe özet 300 kelimeden az olmamak kaydıyla 1-3 sayfa,  İngilizce genişletilmiş özet de 3-5 sayfa arasında olmalıdır. İngilizce tezlerde ise, İngilizce özet 300 kelimeden az olmamak kaydıyla 1-3 sayfa,  Türkçe genişletilmiş özet de 3-5 sayfa arasında olmalıdır. }\textbf{Özetlerde tezde ele alınan konu kısaca tanıtılarak, kullanılan yöntemler ve ulaşılan sonuçlar belirtilir}. Özetlerde kaynak, şekil, tablo verilmez. Özetlerin başında, birinci dereceden başlık formatında tezin adı (önce 72, sonra 18 punto aralık bırakılarak ve 1 satır aralıklı olarak) yazılacaktır.  Başlığın altına büyük harflerle sayfa ortalanarak (Türkçe özet için) \textbf{ÖZET} ve (İngilizce özet için) \textbf{SUMMARY} yazılmalıdır. Türkçe tezlerde Türkçe özetin İngilizce özetten önce olması önerilir.

Lorem ipsum dolor sit amet, consetetur sadipscing elitr, sed diam nonumy eirmod tempor invidunt ut labore et dolore magna aliquyam erat, sed diam voluptua. At vero eos et accusam et justo duo dolores et ea rebum. Stet clita kasd gub rgren, no sea takimata sanctus est Lorem ipsum dolor sit amet, consetetur sadipscing elitr, sed diam nonumy eirmod tempor invidunt ut lab ore sit et dolore magna.

Lorem ipsum dolor sit amet, consetetur sadipscing elitr, sed diam nonumy eirmod tempor invidunt ut labore et dolore magna aliquyam erat, sed diam voluptua. At vero eos et accusam et justo duo dolores et ea rebum. Stet clita kasd gub rgren, no sea takimata sanctus est Lorem ipsum dolor sit amet, consetetur sadipscing elitr, sed diam nonumy eirmod tempor invidunt ut lab ore sit et dolore magna.

Lorem ipsum dolor sit amet, consetetur sadipscing elitr, sed diam nonumy eirmod tempor invidunt ut labore et dolore magna aliquyam erat, sed diam voluptua. At vero eos et accusam et justo duo dolores et ea rebum. Stet clita kasd gub rgren, no sea takimata sanctus est Lorem ipsum dolor sit amet, consetetur sadipscing elitr, sed diam nonumy eirmod tempor invidunt ut lab ore sit et dolore magna.

Lorem ipsum dolor sit amet, consetetur sadipscing elitr, sed diam nonumy eirmod tempor invidunt ut labore et dolore magna aliquyam erat, sed diam voluptua. At vero eos et accusam et justo duo dolores et ea rebum. Stet clita kasd gub rgren, no sea takimata sanctus est Lorem ipsum dolor sit amet, consetetur sadipscing elitr, sed diam nonumy eirmod tempor invidunt ut lab ore sit et dolore magna.

Lorem ipsum dolor sit amet, consetetur sadipscing elitr, sed diam nonumy eirmod tempor invidunt ut labore et dolore magna aliquyam erat, sed diam voluptua. At vero eos et accusam et justo duo dolores et ea rebum. Stet clita kasd gub rgren, no sea takimata sanctus est Lorem ipsum dolor sit amet, consetetur sadipscing elitr, sed diam nonumy eirmod tempor invidunt ut lab ore sit et dolore magna.

Lorem ipsum dolor sit amet, consetetur sadipscing elitr, sed diam nonumy eirmod tempor invidunt ut labore et dolore magna aliquyam erat, sed diam voluptua. At vero eos et accusam et justo duo dolores et ea rebum. Stet clita kasd gub rgren, no sea takimata sanctus est Lorem ipsum dolor sit amet, consetetur sadipscing elitr, sed diam nonumy eirmod tempor invidunt ut lab ore sit et dolore magna.

Lorem ipsum dolor sit amet, consetetur sadipscing elitr, sed diam nonumy eirmod tempor invidunt ut labore et dolore magna aliquyam erat, sed diam voluptua. At vero eos et accusam et justo duo dolores et ea rebum. Stet clita kasd gub rgren, no sea takimata sanctus est Lorem ipsum dolor sit amet, consetetur sadipscing elitr, sed diam nonumy eirmod tempor invidunt ut lab ore sit et dolore magna.

Lorem ipsum dolor sit amet, consetetur sadipscing elitr, sed diam nonumy eirmod tempor invidunt ut labore et dolore magna aliquyam erat, sed diam voluptua. At vero eos et accusam et justo duo dolores et ea rebum. Stet clita kasd gub rgren, no sea takimata sanctus est Lorem ipsum dolor sit amet, consetetur sadipscing elitr, sed diam nonumy eirmod tempor invidunt ut lab ore sit et dolore magna.

Lorem ipsum dolor sit amet, consetetur sadipscing elitr, sed diam nonumy eirmod tempor invidunt ut labore et dolore magna aliquyam erat, sed diam voluptua. At vero eos et accusam et justo duo dolores et ea rebum. Stet clita kasd gub rgren, no sea takimata sanctus est Lorem ipsum dolor sit amet, consetetur sadipscing elitr, sed diam nonumy eirmod tempor invidunt ut lab ore sit et dolore magna.

Lorem ipsum dolor sit amet, consetetur sadipscing elitr, sed diam nonumy eirmod tempor invidunt ut labore et dolore magna aliquyam erat, sed diam voluptua. At vero eos et accusam et justo duo dolores et ea rebum. Stet clita kasd gub rgren, no sea takimata sanctus est Lorem ipsum dolor sit amet, consetetur sadipscing elitr, sed diam nonumy eirmod tempor invidunt ut lab ore sit et dolore magna.

Lorem ipsum dolor sit amet, consetetur sadipscing elitr, sed diam nonumy eirmod tempor invidunt ut labore et dolore magna aliquyam erat, sed diam voluptua. At vero eos et accusam et justo duo dolores et ea rebum. Stet clita kasd gub rgren, no sea takimata sanctus est Lorem ipsum dolor sit amet, consetetur sadipscing elitr, sed diam nonumy eirmod tempor invidunt ut lab ore sit et dolore magna.

Lorem ipsum dolor sit amet, consetetur sadipscing elitr, sed diam nonumy eirmod tempor invidunt ut labore et dolore magna aliquyam erat, sed diam voluptua. At vero eos et accusam et justo duo dolores et ea rebum. Stet clita kasd gub rgren, no sea takimata sanctus est Lorem ipsum dolor sit amet, consetetur sadipscing elitr, sed diam nonumy eirmod tempor invidunt ut lab ore sit et dolore magna.

Lorem ipsum dolor sit amet, consetetur sadipscing elitr, sed diam nonumy eirmod tempor invidunt ut labore et dolore magna aliquyam erat, sed diam voluptua. At vero eos et accusam et justo duo dolores et ea rebum. Stet clita kasd gub rgren, no sea takimata sanctus est Lorem ipsum dolor sit amet, consetetur sadipscing elitr, sed diam nonumy eirmod tempor invidunt ut lab ore sit et dolore magna.

Lorem ipsum dolor sit amet, consetetur sadipscing elitr, sed diam nonumy eirmod tempor invidunt ut labore et dolore magna aliquyam erat, sed diam voluptua. At vero eos et accusam et justo duo dolores et ea rebum. Stet clita kasd gub rgren, no sea takimata sanctus est Lorem ipsum dolor sit amet, consetetur sadipscing elitr, sed diam nonumy eirmod tempor invidunt ut lab ore sit et dolore magna.

Lorem ipsum dolor sit amet, consetetur sadipscing elitr, sed diam nonumy eirmod tempor invidunt ut labore et dolore magna aliquyam erat, sed diam voluptua. At vero eos et accusam et justo duo dolores et ea rebum. Stet clita kasd gub rgren, no sea takimata sanctus est Lorem ipsum dolor sit amet, consetetur sadipscing elitr, sed diam nonumy eirmod tempor invidunt ut lab ore sit et dolore magna.

Lorem ipsum dolor sit amet, consetetur sadipscing elitr, sed diam nonumy eirmod tempor invidunt ut labore et dolore magna aliquyam erat, sed diam voluptua. At vero eos et accusam et justo duo dolores et ea rebum. Stet clita kasd gub rgren, no sea takimata sanctus est Lorem ipsum dolor sit amet, consetetur sadipscing elitr, sed diam nonumy eirmod tempor invidunt ut lab ore sit et dolore magna.

Lorem ipsum dolor sit amet, consetetur sadipscing elitr, sed diam nonumy eirmod tempor invidunt ut labore et dolore magna aliquyam erat, sed diam voluptua. At vero eos et accusam et justo duo dolores et ea rebum. Stet clita kasd gub rgren, no sea takimata sanctus est Lorem ipsum dolor sit amet, consetetur sadipscing elitr, sed diam nonumy eirmod tempor invidunt ut lab ore sit et dolore magna.

Lorem ipsum dolor sit amet, consetetur sadipscing elitr, sed diam nonumy eirmod tempor invidunt ut labore et dolore magna aliquyam erat, sed diam voluptua. At vero eos et accusam et justo duo dolores et ea rebum. Stet clita kasd gub rgren, no sea takimata sanctus est Lorem ipsum dolor sit amet, consetetur sadipscing elitr, sed diam nonumy eirmod tempor invidunt ut lab ore sit et dolore magna.

Lorem ipsum dolor sit amet, consetetur sadipscing elitr, sed diam nonumy eirmod tempor invidunt ut labore et dolore magna aliquyam erat, sed diam voluptua. At vero eos et accusam et justo duo dolores et ea rebum. Stet clita kasd gub rgren, no sea takimata sanctus est Lorem ipsum dolor sit amet, consetetur sadipscing elitr, sed diam nonumy eirmod tempor invidunt ut lab ore sit et dolore magna.

Lorem ipsum dolor sit amet, consetetur sadipscing elitr, sed diam nonumy eirmod tempor invidunt ut labore et dolore magna aliquyam erat, sed diam voluptua. At vero eos et accusam et justo duo dolores et ea rebum. Stet clita kasd gub rgren, no sea takimata sanctus est Lorem ipsum dolor sit amet, consetetur sadipscing elitr, sed diam nonumy eirmod tempor invidunt ut lab ore sit et dolore magna.
\end{ozet}

\begin{summary}
\textcolor[HTML]{FF0000}{1-line spacing must be set for summaries. For theses in Turkish, the summary in Turkish must have 400 words minimum and span 1 to 3 pages, whereas the extended summary in English must span 3-5 pages. For theses in English, the summary in English must have 400 words minimum and span 1-3 pages, whereas the extended summary in Turkish must span 3-5 pages. }A summary must briefly mention the subject of the thesis, the method(s) used and the conclusions derived. References, figures and tables must not be given in Summary. Above the Summary, the thesis title in first level title format (i.e., 72 pt before and 18 pt after paragraph spacing, and 1 line spacing) must be placed. Below the title, the expression \textbf{ÖZET }(for summary in Turkish) and \textbf{SUMMARY }(for summary in English) must be written horizontally centered. It is recommended that the summary in English is placed before the summary in Turkish.

Lorem ipsum dolor sit amet, consetetur sadipscing elitr, sed diam nonumy eirmod tempor invidunt ut labore et dolore magna aliquyam erat, sed diam voluptua. At vero eos et accusam et justo duo dolores et ea rebum. Stet clita kasd gub rgren, no sea takimata sanctus est Lorem ipsum dolor sit amet, consetetur sadipscing elitr, sed diam nonumy eirmod tempor invidunt ut lab ore sit et dolore magna.

Lorem ipsum dolor sit amet, consetetur sadipscing elitr, sed diam nonumy eirmod tempor invidunt ut labore et dolore magna aliquyam erat, sed diam voluptua. At vero eos et accusam et justo duo dolores et ea rebum. Stet clita kasd gub rgren, no sea takimata sanctus est Lorem ipsum dolor sit amet, consetetur sadipscing elitr, sed diam nonumy eirmod tempor invidunt ut lab ore sit et dolore magna.

Lorem ipsum dolor sit amet, consetetur sadipscing elitr, sed diam nonumy eirmod tempor invidunt ut labore et dolore magna aliquyam erat, sed diam voluptua. At vero eos et accusam et justo duo dolores et ea rebum. Stet clita kasd gub rgren, no sea takimata sanctus est Lorem ipsum dolor sit amet, consetetur sadipscing elitr, sed diam nonumy eirmod tempor invidunt ut lab ore sit et dolore magna.

Lorem ipsum dolor sit amet, consetetur sadipscing elitr, sed diam nonumy eirmod tempor invidunt ut labore et dolore magna aliquyam erat, sed diam voluptua. At vero eos et accusam et justo duo dolores et ea rebum. Stet clita kasd gub rgren, no sea takimata sanctus est Lorem ipsum dolor sit amet, consetetur sadipscing elitr, sed diam nonumy eirmod tempor invidunt ut lab ore sit et dolore magna.

Lorem ipsum dolor sit amet, consetetur sadipscing elitr, sed diam nonumy eirmod tempor invidunt ut labore et dolore magna aliquyam erat, sed diam voluptua. At vero eos et accusam et justo duo dolores et ea rebum. Stet clita kasd gub rgren, no sea takimata sanctus est Lorem ipsum dolor sit amet, consetetur sadipscing elitr, sed diam nonumy eirmod tempor invidunt ut lab ore sit et dolore magna.

Lorem ipsum dolor sit amet, consetetur sadipscing elitr, sed diam nonumy eirmod tempor invidunt ut labore et dolore magna aliquyam erat, sed diam voluptua. At vero eos et accusam et justo duo dolores et ea rebum. Stet clita kasd gub rgren, no sea takimata sanctus est Lorem ipsum dolor sit amet, consetetur sadipscing elitr, sed diam nonumy eirmod tempor invidunt ut lab ore sit et dolore magna.

Lorem ipsum dolor sit amet, consetetur sadipscing elitr, sed diam nonumy eirmod tempor invidunt ut labore et dolore magna aliquyam erat, sed diam voluptua. At vero eos et accusam et justo duo dolores et ea rebum. Stet clita kasd gub rgren, no sea takimata sanctus est Lorem ipsum dolor sit amet, consetetur sadipscing elitr, sed diam nonumy eirmod tempor invidunt ut lab ore sit et dolore magna.

Lorem ipsum dolor sit amet, consetetur sadipscing elitr, sed diam nonumy eirmod tempor invidunt ut labore et dolore magna aliquyam erat, sed diam voluptua. At vero eos et accusam et justo duo dolores et ea rebum. Stet clita kasd gub rgren, no sea takimata sanctus est Lorem ipsum dolor sit amet, consetetur sadipscing elitr, sed diam nonumy eirmod tempor invidunt ut lab ore sit et dolore magna.

Lorem ipsum dolor sit amet, consetetur sadipscing elitr, sed diam nonumy eirmod tempor invidunt ut labore et dolore magna aliquyam erat, sed diam voluptua. At vero eos et accusam et justo duo dolores et ea rebum. Stet clita kasd gub rgren, no sea takimata sanctus est Lorem ipsum dolor sit amet, consetetur sadipscing elitr, sed diam nonumy eirmod tempor invidunt ut lab ore sit et dolore magna.

Lorem ipsum dolor sit amet, consetetur sadipscing elitr, sed diam nonumy eirmod tempor invidunt ut labore et dolore magna aliquyam erat, sed diam voluptua. At vero eos et accusam et justo duo dolores et ea rebum. Stet clita kasd gub rgren, no sea takimata sanctus est Lorem ipsum dolor sit amet, consetetur sadipscing elitr, sed diam nonumy eirmod tempor invidunt ut lab ore sit et dolore magna.

Lorem ipsum dolor sit amet, consetetur sadipscing elitr, sed diam nonumy eirmod tempor invidunt ut labore et dolore magna aliquyam erat, sed diam voluptua. At vero eos et accusam et justo duo dolores et ea rebum. Stet clita kasd gub rgren, no sea takimata sanctus est Lorem ipsum dolor sit amet, consetetur sadipscing elitr, sed diam nonumy eirmod tempor invidunt ut lab ore sit et dolore magna.

Lorem ipsum dolor sit amet, consetetur sadipscing elitr, sed diam nonumy eirmod tempor invidunt ut labore et dolore magna aliquyam erat, sed diam voluptua. At vero eos et accusam et justo duo dolores et ea rebum. Stet clita kasd gub rgren, no sea takimata sanctus est Lorem ipsum dolor sit amet, consetetur sadipscing elitr, sed diam nonumy eirmod tempor invidunt ut lab ore sit et dolore magna.

Lorem ipsum dolor sit amet, consetetur sadipscing elitr, sed diam nonumy eirmod tempor invidunt ut labore et dolore magna aliquyam erat, sed diam voluptua. At vero eos et accusam et justo duo dolores et ea rebum. Stet clita kasd gub rgren, no sea takimata sanctus est Lorem ipsum dolor sit amet, consetetur sadipscing elitr, sed diam nonumy eirmod tempor invidunt ut lab ore sit et dolore magna.

Lorem ipsum dolor sit amet, consetetur sadipscing elitr, sed diam nonumy eirmod tempor invidunt ut labore et dolore magna aliquyam erat, sed diam voluptua. At vero eos et accusam et justo duo dolores et ea rebum. Stet clita kasd gub rgren, no sea takimata sanctus est Lorem ipsum dolor sit amet, consetetur sadipscing elitr, sed diam nonumy eirmod tempor invidunt ut lab ore sit et dolore magna.

Lorem ipsum dolor sit amet, consetetur sadipscing elitr, sed diam nonumy eirmod tempor invidunt ut labore et dolore magna aliquyam erat, sed diam voluptua. At vero eos et accusam et justo duo dolores et ea rebum. Stet clita kasd gub rgren, no sea takimata sanctus est Lorem ipsum dolor sit amet, consetetur sadipscing elitr, sed diam nonumy eirmod tempor invidunt ut lab ore sit et dolore magna.

Lorem ipsum dolor sit amet, consetetur sadipscing elitr, sed diam nonumy eirmod tempor invidunt ut labore et dolore magna aliquyam erat, sed diam voluptua. At vero eos et accusam et justo duo dolores et ea rebum. Stet clita kasd gub rgren, no sea takimata sanctus est Lorem ipsum dolor sit amet, consetetur sadipscing elitr, sed diam nonumy eirmod tempor invidunt ut lab ore sit et dolore magna.

Lorem ipsum dolor sit amet, consetetur sadipscing elitr, sed diam nonumy eirmod tempor invidunt ut labore et dolore magna aliquyam erat, sed diam voluptua. At vero eos et accusam et justo duo dolores et ea rebum. Stet clita kasd gub rgren, no sea takimata sanctus est Lorem ipsum dolor sit amet, consetetur sadipscing elitr, sed diam nonumy eirmod tempor invidunt ut lab ore sit et dolore magna.

Lorem ipsum dolor sit amet, consetetur sadipscing elitr, sed diam nonumy eirmod tempor invidunt ut labore et dolore magna aliquyam erat, sed diam voluptua. At vero eos et accusam et justo duo dolores et ea rebum. Stet clita kasd gub rgren, no sea takimata sanctus est Lorem ipsum dolor sit amet, consetetur sadipscing elitr, sed diam nonumy eirmod tempor invidunt ut lab ore sit et dolore magna.

Lorem ipsum dolor sit amet, consetetur sadipscing elitr, sed diam nonumy eirmod tempor invidunt ut labore et dolore magna aliquyam erat, sed diam voluptua. At vero eos et accusam et justo duo dolores et ea rebum. Stet clita kasd gub rgren, no sea takimata sanctus est Lorem ipsum dolor sit amet, consetetur sadipscing elitr, sed diam nonumy eirmod tempor invidunt ut lab ore sit et dolore magna.

Lorem ipsum dolor sit amet, consetetur sadipscing elitr, sed diam nonumy eirmod tempor invidunt ut labore et dolore magna aliquyam erat, sed diam voluptua. At vero eos et accusam et justo duo dolores et ea rebum. Stet clita kasd gub rgren, no sea takimata sanctus est Lorem ipsum dolor sit amet, consetetur sadipscing elitr, sed diam nonumy eirmod tempor invidunt ut lab ore sit et dolore magna.

Lorem ipsum dolor sit amet, consetetur sadipscing elitr, sed diam nonumy eirmod tempor invidunt ut labore et dolore magna aliquyam erat, sed diam voluptua. At vero eos et accusam et justo duo dolores et ea rebum. Stet clita kasd gub rgren, no sea takimata sanctus est Lorem ipsum dolor sit amet, consetetur sadipscing elitr, sed diam nonumy eirmod tempor invidunt ut lab ore sit et dolore magna.

Lorem ipsum dolor sit amet, consetetur sadipscing elitr, sed diam nonumy eirmod tempor invidunt ut labore et dolore magna aliquyam erat, sed diam voluptua. At vero eos et accusam et justo duo dolores et ea rebum. Stet clita kasd gub rgren, no sea takimata sanctus est Lorem ipsum dolor sit amet, consetetur sadipscing elitr, sed diam nonumy eirmod tempor invidunt ut lab ore sit et dolore magna.
\end{summary}


%----------------------------------------------------------------- ana metin
\anametin
\chapter{GİRİŞ (BİRİNCİ DERECE BAŞLIK)}
Birinci dereceden başlıklar okuma yönünde, sağ sayfadan başlamalı, büyük ve koyu harflerle yazılmalıdır. (Örnek: \textbf{1. GİRİŞ})

Lorem ipsum dolor sit amet, consetetur sadipscing elitr, sed diam nonumy eirmod tempor invidunt ut labore et dolore magna aliquyam erat, sed diam voluptua. At vero eos et accusam et justo duo dolores et ea rebum. Stet clita kasd gub rgren, no sea takimata sanctus est Lorem ipsum dolor sit amet, consetetur sadipscing elitr, sed diam nonumy eirmod tempor invidunt ut lab ore sit et dolore magna.

\section{Tezin Kapsamı (İkinci Derece Başlık Nasıl: İlk Harfler Büyük, İkinci Derece Başlık Nasıl: İlk Harfler Büyük)}
Lorem ipsum dolor sit amet, consetetur sadipscing elitr, sed diam nonumy eirmod tempor invidunt ut labore et dolore magna aliquyam erat, sed diam voluptua. At vero eos et accusam et justo duo dolores et ea rebum. Stet clita kasd gub rgren, no sea takimata sanctus est Lorem ipsum dolor sit amet, consetetur sadipscing elitr, sed diam nonumy eirmod tempor invidunt ut lab ore sit et dolore magna.

\section{Tezin Amacı}
Lorem ipsum dolor sit amet, consetetur sadipscing elitr, sed diam nonumy eirmod tempor invidunt ut labore et dolore magna aliquyam erat, sed diam voluptua. At vero eos et accusam et justo duo dolores et ea rebum. Stet clita kasd gub rgren, no sea takimata sanctus est.

Lorem ipsum dolor sit amet, consetetur sadipscing elitr, sed diam nonumy eirmod tempor invidunt ut labore et dolore magna aliquyam erat, sed diam voluptua. At vero eos et accusam et justo duo dolores et ea rebum. Stet clita kasd gub rgren, no sea takimata sanctus est Lorem ipsum dolor sit amet, consetetur sadipscing elitr, sed diam nonumy eirmod tempor invidunt ut lab ore sit et dolore magna. Lorem ipsum dolor sit amet, consetetur sadipscing elitr, sed diam nonumy eirmod tempor invidunt.

\subsection{Tezin ikincil amacı}
Lorem ipsum dolor sit amet, consetetur sadipscing elitr, sed diam nonumy eirmod tempor invidunt ut labore et dolore magna aliquyam erat, sed diam voluptua. At vero eos et accusam et justo duo dolores et ea rebum. Stet clita kasd gub rgren, no sea takimata sanctus est.

\subsubsection{Dördüncü derece başlık nasıl: ilk harf büyük diğerleri küçük}
Lorem ipsum dolor sit amet, consetetur sadipscing elitr, sed diam nonumy eirmod tempor invidunt ut lab ore sit et dolore magna. Stet clita kasd gub rgren, no sea takimata sanctus est Lorem ipsum dolor sit amet, consetetur sadipscing elitr, sed diam nonumy eirmod tempor invidunt ut lab ore sit et dolore magna. Stet clita kasd gub rgren, no sea takimata sanctus est.

\subsubsection{Dördüncü derece başlık nasıl: ilk harf büyük diğerleri küçük}
Lorem ipsum dolor sit amet, consetetur sadipscing elitr, sed diam nonumy eirmod tempor invidunt ut lab ore sit et dolore magna. Stet clita kasd gub rgren, no sea takimata sanctus est Lorem ipsum dolor sit amet, consetetur sadipscing elitr, sed diam nonumy eirmod tempor invidunt ut lab ore sit et dolore magna. Stet clita kasd gub rgren, no sea takimata sanctus est. Lorem ipsum dolor sit amet, consetetur sadipscing elitr, sed diam nonumy eirmod tempor invidunt ut lab ore sit et dolore magna. Stet clita kasd gub rgren, no sea takimata sanctus est. Lorem ipsum dolor sit amet, consetetur sadipscing elitr, sed diam nonumy eirmod tempor invidunt ut lab ore sit et dolore magna. Stet clita kasd gub rgren, no sea takimata sanctus est. Lorem ipsum dolor sit amet, consetetur sadipscing elitr, sed diam nonumy.

\baslikbes{Beşinci derece başlık nasıl: ilk harf büyük diğerleri küçük}
Lorem ipsum dolor sit amet, consetetur sadipscing elitr, sed diam nonumy eirmod tempor invidunt ut lab ore sit et dolore magna. Stet clita kasd gub rgren, no sea takimata sanctus est Lorem ipsum dolor sit amet, consetetur sadipscing elitr, sed diam nonumy eirmod tempor invidunt ut lab ore sit et dolore magna. Stet clita kasd gub rgren, no sea takimata sanctus est. Lorem ipsum dolor sit amet, consetetur sadipscing elitr, sed diam nonumy eirmod tempor invidunt ut lab ore sit et dolore magna. Stet clita kasd gub rgren, no sea takimata sanctus est. Lorem ipsum dolor sit amet, consetetur sadipscing elitr, sed diam nonumy eirmod tempor invidunt ut lab ore sit et dolore magna. Stet clita kasd gub rgren, no sea takimata sanctus est. Lorem ipsum dolor sit amet, consetetur sadipscing elitr, sed diam nonumy.

Lorem ipsum dolor sit amet, consetetur sadipscing elitr, sed diam nonumy eirmod tempor invidunt ut lab ore sit et dolore magna. Stet clita kasd gub rgren, no sea takimata sanctus est Lorem ipsum dolor sit amet, consetetur sadipscing elitr, sed diam nonumy eirmod tempor invidunt ut lab ore sit et dolore magna. Stet clita kasd gub rgren, no sea takimata sanctus est.

\section{Literatür Araştırması}
Lorem ipsum dolor sit amet, consetetur sadipscing elitr, sed diam nonumy eirmod tempor invidunt ut labore et dolore magna aliquyam erat, sed diam voluptua. At vero eos et accusam et justo duo dolores et ea rebum. Stet clita kasd gub rgren, no sea takimata sanctus est Lorem ipsum dolor sit amet, consetetur sadipscing elitr, sed diam nonumy eirmod tempor invidunt ut labore et dolore magna aliquyam erat, sed diam voluptua. At vero eos et accusam et justo duo dolores et ea rebum. Stet clita kasd gub rgren, no sea takimata sanctus est Lorem ipsum dolor sit amet, consetetur sadipscing elitr, sed diam nonumy eirmod tempor invidunt ut labore et dolore magna aliquyam erat, sed diam voluptua. At vero eos et accusam et justo duo dolores et ea rebum. Stet clita kasd gub rgren, no sea takimata sanctus est Lorem ipsum dolor sit amet, consetetur sadipscing elitr, sed diam nonumy eirmod tempor invidunt ut labore et dolore magna aliquyam erat, sed diam voluptua. At vero eos et accusam et justo duo dolores et ea rebum. Stet clita kasd gub rgren, no sea takimata sanctus est.

At vero eos et accusam et justo duo dolores et ea rebum. Stet clita kasd gub rgren, no sea takimata sanctus est Lorem ipsum dolor sit amet, consetetur sadipscing elitr, sed diam nonumy eirmod tempor invidunt ut labore et dolore magna aliquyam erat, sed diam voluptua. At vero eos et accusam et justo duo dolores et ea rebum. Stet clita kasd gub rgren, no sea takimata sanctus est Lorem ipsum dolor sit amet, consetetur sadipscing elitr, sed diam nonumy eirmod tempor invidunt ut labore et dolore magna aliquyam erat, sed diam voluptua. At vero eos et accusam et justo duo dolores et ea rebum. Stet clita kasd gub rgren, no sea takimata sanctus est.

\section{Hipotez}
Lorem ipsum dolor sit amet, consetetur sadipscing elitr, sed diam nonumy eirmod tempor invidunt ut labore et dolore magna aliquyam erat, sed diam voluptua. At vero eos et accusam et justo duo dolores et ea rebum. Stet clita kasd gub rgren, no sea takimata sanctus est Lorem ipsum dolor sit amet, consetetur sadipscing elitr, sed diam nonumy eirmod tempor invidunt ut labore et dolore magna aliquyam erat, sed diam voluptua. At vero eos et accusam et justo duo dolores et ea rebum. Stet clita kasd gub rgren, no sea takimata sanctus est.

Lorem ipsum dolor sit amet, consetetur sadipscing elitr, sed diam nonumy eirmod tempor invidunt ut labore et dolore magna aliquyam erat, sed diam voluptua. At vero eos et accusam et justo duo dolores et ea rebum. Stet clita kasd gub rgren, no sea takimata sanctus est Lorem ipsum dolor sit amet, consetetur sadipscing elitr, sed diam nonumy eirmod tempor invidunt ut labore et dolore magna aliquyam erat, sed diam voluptua. At vero eos et accusam et justo duo dolores et ea rebum. Stet clita kasd gub rgren, no sea takimata sanctus est.


\chapter{ŞEKİL VE TABLOLAR}
\section{Şekil Atıflar ve Şekil Örneği}
Tablo ve Şekiller sayfa düzeni esaslarına uymak koşuluyla, metinde ilk söz edildikleri yerden hemen sonraya veya mümkün olduğu ölçüde en yakınına yerleştirilmelidir. Tablo ve şekillerden önce, ilgili tablo ya da şekile atıfta bulunulmalıdır.

Şekil başlığı şeklin altına, tablo başlığı ise tablonun üstüne yazılır. Tablo ve şekil başlıkları bir satıra sığıyorsa yazı bloğuna göre ortalı olarak yazılır. Bir satırdan uzun olan başlıklar ise iki yana yaslı olarak yazılmalıdır. Bu durumda, ikinci satır, tablo ve şekil numarasından sonra ilk satır metni ile aynı hizadan başlayacak şekilde girintili olmalıdır.

Tablo ve şekillere, ilk rakam bölüm numarası (Ekler bölümünde ise harf), ikinci rakam tablo veya şeklin bölüm içindeki sıra numarası olmak üzere numara verilir. Tablo, şekil kelimeleri ve numaraları koyu harflerle yazılır ve sonlarına nokta konulmalıdır.  (Örnek: Tablo 1.2., Şekil 3.5., Tablo A.1., Şekil B.5.).

Tablo ve şekillere metin içerisinde atıf yapılırken tablo veya şekil numaralarının sonundaki nokta kullanılmaz. (Örneğin; Şekil 1.1'de görüldüğü gibi \ldots{})

Tablo başlıkları, bir satır aralıklı olarak, başlıktan önce 12 punto (nk), sonra ise 6 punto (nk) aralık bırakılarak yazılmalı ve Tablo başlıkları nokta ile bitirilmelidir. Tablo başlığı ile tablo farklı sayfalarda yer alamaz.

Tablodan sonra gelen metin bölümündeki ilk paragraf üstten 12 punto (nk) ve alttan 6 punto (nk) aralık bırakılarak yazılmalıdır. Tablolardan hemen sonra gelecek başlıklar, belirtilen başlık formatlarında değişiklik yapılmadan aynen kullanılmalıdır.

Tablolarda kullanılacak dipnotlar (gerekliyse) 1 satır aralıklı ve metinden 2 yazı boyutu küçük olarak yazılmalıdır.

Şekil başlıkları, bir satır aralıklı olarak, başlıktan önce 6 punto (nk), sonra ise 12 punto (nk) aralık bırakılarak yazılmalı ve şekil başlıkları nokta ile bitirilmelidir. Şekil başlığı ile şekil farklı sayfalarda yer alamaz. Şekilden önce gelen metin bölümündeki son paragraf, önce 6 punto (nk), sonra 12 punto (nk) aralık bırakılarak yazılmalıdır.

Bir sayfayı aşan büyüklükteki tablo ve şekillerde 2. sayfada aynı tablo/şekil numarası ve açıklaması yazılarak, tablo/şekil numarası ile açıklaması arasına, parantez içinde (Devamı) yazılmalıdır. (Örneğin; Tablo 1.1. (Devamı):\textbf{ }Geometrik yapı parametreleri, Şekil 1.1. (Devamı):\textbf{ }Elektronik soğurma spektrumu).

Tezlerde kullanılacak olan grafik ve resimler şekil olarak dikkate alınmalı ve buna göre düzenlenmelidir.

Birden fazla Tablo veya Şekil aynı sayfada yer alabilir. Ancak, 4 sayfadan daha uzun olan tablo veya şekiller ek olarak verilmelidir.

% [Normal] sp=after:240
Katlı sayfa (A4 ebadından büyük olan) ve sayfa üzerine iliştirilmiş görsel malzeme gibi sayfa kalınlığını arttırarak tezin açılma düzenini bozan sayfalar da ek olarak verilmelidir.

\begin{figure}[H]
  \centering
  \includegraphics[width=11.11cm]{sekiller/image1.png}
  \caption{Tüm şekil ve tablolar ile bunların başlıkları yazı bloğuna göre ortalı olarak yerleştirilmelidir.}
\end{figure}
Tablo ve şekillerde gerekli ise yazı boyutu 8 punto yazı boyutuna kadar küçültülebilir.

Tablolar tezde kullanılan yazı karakteriyle hazırlanır, şekillerde kullanılan yazı karakteri tez boyunca kendi içerisinde tutarlı olmalıdır.

% [Normal] sp=after:240
Tablolar ve şekiller sayfa düzeni esaslarına uymak şartı ile metinde ilk söz edildikleri yerden hemen sonraya mümkün olduğu kadar yakın yerleştirilmelidir (Şekil 2.1.). Tablo ve şekillerden önce, ilgili tablo ya da şekile atıfta bulunulmalıdır (Tablo 1.1).

\begin{figure}[H]
  \centering
  \includegraphics[width=11.11cm]{sekiller/image1.png}
  \caption{Bir satıra sığan şekil başlıkları ortalı olarak yerleştirilmelidir.}
\end{figure}
Şekil 2.2’de ki Sed diam nonumy eirmod tempor invidunt ut labore et dolore magna aliquyam erat, sed diam voluptua. At vero eos et accusam et justo duo dolores et ea rebum. Lorem ipsum dolor sit amet, consetetur sadipscing elitr, sed diam nonumy eirmod tempor invidunt ut labore et dolore magna aliquyam erat, sed diam voluptua. At vero eos et accusam et justo duo dolores et ea rebum. At vero eos et accusam et justo duo dolores et ea rebum. At vero eos et accusam et justo duo dolores et ea rebum.

Lorem ipsum dolor sit amet, consetetur sadipscing elitr, sed diam nonumy eirmod tempor invidunt ut labore et dolore magna aliquyam erat, sed diam voluptua. At vero eos et accusam et justo duo dolores et ea rebum. At vero eos et accusam et justo duo dolores et ea rebum. At vero eos et accusam et justo duo dolores et ea rebum.

\section{Yatay Sayfada Şekil Örneği}
Lorem ipsum dolor sit amet, consetetur sadipscing elitr, sed diam nonumy eirmod tempor invidunt ut labore et dolore magna aliquyam erat, sed diam voluptua. At vero eos et accusam et justo duo dolores et ea rebum (Şekil 2.3). Lorem ipsum dolor sit amet, consetetur sadipscing elitr, sed diam nonumy eirmod tempor invidunt ut labore et dolore magna aliquyam erat, sed diam voluptua. At vero eos et accusam et justo duo dolores et ea rebum.

Lorem ipsum dolor sit amet, consetetur sadipscing elitr, sed diam nonumy eirmod tempor invidunt ut labore et dolore magna aliquyam erat, sed diam voluptua. At vero eos et accusam et justo duo dolores et ea rebum. Lorem ipsum dolor sit amet, consetetur sadipscing elitr, sed diam nonumy eirmod tempor invidunt ut labore et dolore magna aliquyam erat, sed diam voluptua. At vero eos et accusam et justo duo dolores et ea rebum. Lorem ipsum dolor sit amet, consetetur sadipscing elitr, sed diam nonumy eirmod tempor invidunt ut labore et dolore magna aliquyam erat, sed diam voluptua. At vero eos et accusam et justo duo dolores et ea rebum.

\begin{landscape}
\begin{figure}[H]
  \centering
  \includegraphics[width=21.96cm,height=10.54cm]{sekiller/image2.png}
  \caption{Yatay tam sayfa şekil.}
\end{figure}
\end{landscape}
Lorem ipsum dolor sit amet, consetetur sadipscing elitr, sed diam nonumy eirmod tempor invidunt ut labore et dolore magna aliquyam erat, sed diam voluptua. At vero eos et accusam et justo duo dolores et ea rebum. Lorem ipsum dolor sit amet, consetetur sadipscing elitr, sed diam nonumy eirmod tempor invidunt ut labore et dolore magna aliquyam erat, sed diam voluptua. At vero eos et accusam et justo duo dolores et ea rebum. Lorem ipsum dolor sit amet, consetetur sadipscing elitr, sed diam nonumy eirmod tempor invidunt ut labore et dolore magna aliquyam erat, sed diam voluptua. At vero eos et accusam et justo duo dolores et ea rebum. Lorem ipsum dolor sit amet, consetetur sadipscing elitr, sed diam nonumy eirmod tempor invidunt ut labore et dolore magna aliquyam erat, sed diam voluptua. At vero eos et accusam et justo duo dolores et ea rebum.

\section{Tablo Atıfları ve Tablo Örneği}
Lorem ipsum dolor sit amet, consetetur sadipscing elitr, sed diam nonumy eirmod tempor invidunt ut labore et dolore magna aliquyam erat, sed diam voluptua. At vero eos et accusam et justo duo dolores et ea rebum. Stet clita kasd gub rgren, no sea takimata sanctus est Lorem ipsum dolor sit amet, consetetur sadipscing elitr, sed diam nonumy eirmod tempor invidunt ut lab ore sit et dolore magna.

Tablolar sayfa düzeni esaslarına uymak koşuluyla metinde ilk söz edildikleri yerden hemen sonraya veya mümkün olduğu ölçüde en yakınına yerleştirilmelidir. Tablolardan önce, ilgili tabloya atıfta bulunulmalıdır.

Her tablonun başlığı tablonun üstüne yazılır. Tüm Tablo başlıkları bir satıra sığıyorsa yazı bloğuna göre ortalı olarak yerleştirilmelidir. Bir satırdan uzun olan başlıklar ise iki yana yaslı olarak yerleştirilmelidir. İkinci satır, tablo ve şekil numarasından sonra ilk satır metni ile aynı hizadan başlayacak şekilde girintili olmalıdır.

Tablolarda ilk rakam bölüm numarası (Ekler bölümünde harf), ikinci rakam tabloların  bölüm içindeki sıra numarası olmak üzere numara verilir. Örneklerde olduğu gibi tablo kelimesi ve numaraları koyu harflerle yazılır ve sonunda nokta konulmalıdır.  (Örnek: Tablo 1.2., Tablo A.1.).

Tablo başlıkları, bir satır aralığı kullanılarak ve yazıdan önce 12 punto (nk), sonra 6 punto (nk)  aralık bırakılarak yazılmalı ve Tablo başlığı nokta ile bitirilmelidir. Tablo başlığı ile tablonun tamamı aynı sayfa içinde yer almalıdır.

Tablodan sonra gelen metin bölümündeki ilk paragraf üstten 12 punto (nk) ve alttan 6 punto (nk) aralık bırakılarak yazılmalıdır. Tablolardan hemen sonra gelecek başlıklar, belirtilen başlık formatlarında değişiklik yapılmadan aynen kullanılmalıdır.

Tablolarda dipnot kullanılması gerekiyorsa 1 satır aralıklı ve metinden 2 yazı boyutu küçük yazılmalıdır.

Bir sayfayı aşan büyüklükteki tablolarda 2. sayfada aynı tablo numarası ve açıklaması yazılarak, tablo numarası ile açıklaması arasına, parantez içinde (Devamı) yazılmalıdır. (Örneğin; Tablo 1.1. (Devamı)\textbf{ }Atıklardaki metal içerikleri).

Birden fazla Tablo aynı sayfada yer alabilir. Ancak 4 sayfadan daha uzun olan tablolar ek olarak verilmelidir.

Katlı sayfa (A4 ebadından büyük olan) ve sayfa üzerine iliştirilmiş görsel malzeme gibi sayfa kalınlığını arttırarak tezin açılma düzenini bozan sayfalar ek olarak verilmelidir.

Tablo 2.1’de görüldüğü üzere Lorem ipsum dolor sit amet, consetetur sadipscing elitr, sed diam nonumy eirmod tempor invidunt ut labore et dolore magna aliquyam erat, sed diam voluptua. At vero eos et accusam et justo duo dolores et ea rebum. Stet clita kasd gub rgren, no sea takimata sanctus est Lorem ipsum dolor sit amet, consetetur sadipscing elitr, sed diam nonumy eirmod tempor invidunt ut lab ore sit et dolore magna.

\begin{table}[H]
  \caption{Tek satırlı ve kolonlar ortalanmış tablo.}
  \centering
\begin{tabular}{C{3.09cm}C{2.13cm}C{2.19cm}C{2.56cm}}
\hline\hline
\Ns Kolon A & Kolon B & Kolon C & Kolon D \\
\hline
\Ns Satır A & Satır A & Satır A & Satır A \\
\Ns Satır B & Satır B & Satır B & Satır B \\
\Ns Satır C & Satır C & Satır C & Satır C \\
\hline
\end{tabular}
\end{table}
% [Normal] sp=before:240
Lorem ipsum dolor sit amet, consetetur sadipscing elitr, sed diam nonumy eirmod tempor invidunt ut labore et dolore magna aliquyam erat, sed diam voluptua. At vero eos et accusam et justo duo dolores et ea rebum.

Lorem ipsum dolor sit amet, consetetur sadipscing elitr, sed diam nonumy eirmod tempor invidunt ut labore et dolore magna aliquyam erat, sed diam voluptua. At vero eos et accusam et justo duo dolores et ea rebum. Stet clita kasd gub rgren, no sea takimata sanctus est Lorem ipsum dolor sit amet, consetetur sadipscing elitr, sed diam nonumy eirmod tempor invidunt ut lab ore sit et dolore magna. Lorem ipsum dolor sit amet, consetetur sadipscing elitr, sed diam nonumy eirmod tempor invidunt ut labore et dolore magna aliquyam erat, sed diam voluptua. At vero eos et accusam et justo duo dolores et ea rebum.

\begin{table}[H]
  \caption{Tablo başlığı nokta ile bitirilmelidir.}
  \centering
\begin{tabular}{C{3.05cm}C{2.10cm}C{2.16cm}C{2.52cm}}
\hline\hline
\Ns Kolon A & Kolon B & Kolon C & Kolon D \\
\hline
\Ns Satır A & Satır A & Satır A & Satır A \\
\Ns Satır B & Satır B & Satır B & Satır B \\
\Ns Satır C & Satır C & Satır C & Satır C \\
\hline
\end{tabular}
\end{table}
% [GOVDE] sp=before:240
Lorem ipsum dolor sit amet, consetetur sadipscing elitr, sed diam nonumy eirmod tempor invidunt ut labore et dolore magna aliquyam erat, sed diam voluptua. At vero eos et accusam et justo duo dolores et ea rebum. Lorem ipsum dolor sit amet, consetetur sadipscing elitr, sed diam nonumy eirmod tempor invidunt ut labore et dolore magna aliquyam erat, sed diam voluptua. At vero eos et accusam et justo duo dolores et ea rebum. Lorem ipsum dolor sit amet, consetetur sadipscing elitr, sed diam nonumy eirmod tempor invidunt ut labore et dolore magna aliquyam erat, sed diam voluptua. At vero eos et accusam et justo duo dolores et ea rebum.

\begin{table}[H]
  \caption{Düşey sayfada devam eden tablo örneği.}
  \centering
\begin{tabular}{L{1.72cm}C{1.80cm}C{2.22cm}R{1.04cm}R{1.04cm}R{1.05cm}}
\hline\hline
\multirow{2}{*}{Parametre} & \multirow{2}{*}{Kolon 2} & \multirow{2}{*}{Kolon 3} & \multicolumn{3}{c}{\Ts Kolon 4} \\
\cline{4-6}
 &  &  & \Ts Alt kolon & Alt kolon & Alt kolon \\
\hline
\Ts Satır 1 & -7.680442 & 7.6986348 & 0.00 & 0.00 & 0.00 \\
\Ts Satır 2 & 140 & - & 0.50 & 0.00 & 0.00 \\
\Ts Satır 3 & 37.174357 & 37.16192697 & 0.00 & 0.00 & 0.00 \\
\Ts Satır 4 & 140 & - & 0.50 & 0.00 & 0.00 \\
\Ts Satır 5 & 37.174357 & 37.16192697 & 0.00 & 0.00 & 0.00 \\
\hline
\end{tabular}
\end{table}
\begin{table}[H]
  \ContinuedFloat
  \caption[]{\devami{} Düşey sayfada devam eden tablo örneği.}
  \centering
\begin{tabular}{L{1.72cm}C{1.80cm}C{2.22cm}R{1.04cm}R{1.04cm}R{1.05cm}}
\hline\hline
\multirow{2}{*}{Parametre} & \multirow{2}{*}{Kolon 2} & \multirow{2}{*}{Kolon 3} & \multicolumn{3}{c}{\Ts Kolon 4} \\
\cline{4-6}
 &  &  & \Ts Alt kolon & Alt kolon & Alt kolon \\
\hline
\Ts Satır 1 & -7.680442 & 7.6986348 & 0.00 & 0.00 & 0.00 \\
\Ts Satır 2 & 140 & - & 0.50 & 0.00 & 0.00 \\
\Ts Satır 3 & 37.174357 & 37.16192697 & 0.00 & 0.00 & 0.00 \\
\Ts Satır 4 & 140 & - & 0.50 & 0.00 & 0.00 \\
\Ts Satır 5 & 37.174357 & 37.16192697 & 0.00 & 0.00 & 0.00 \\
\Ts Satır 6 & 140 & - & 0.50 & 0.00 & 0.00 \\
\Ts Satır 7 & 37.174357 & 37.16192697 & 0.00 & 0.00 & 0.00 \\
\hline
\end{tabular}
\end{table}
% [GOVDE] sp=before:240
Loremipsum dolor sit amet, consetetur sadipscing elitr, sed diam nonumy eirmod tempor invidunt ut labore et dolore magna aliquyam erat, sed diam voluptua. At vero eos et accusam et justo duo dolores et ea rebum. Lorem ipsum dolor sit amet, consetetur sadipscing elitr, sed diam nonumy eirmod tempor invidunt ut labore et dolore magna aliquyam erat, sed diam voluptua. At vero eos et accusam et justo duo dolores et ea rebum. Lorem ipsum dolor sit amet, consetetur sadipscing elitr, sed diam nonumy eirmod tempor invidunt ut labore et dolore magna aliquyam erat, sed diam voluptua. At vero eos et accusam et justo duo dolores et ea rebum.

\section{Yatay Sayfada Tablo Örneği}
% [Normal] sp=before:240
Lorem ipsum dolor sit amet, consetetur sadipscing elitr, sed diam nonumy eirmod tempor invidunt ut labore et dolore magna aliquyam erat, sed diam voluptua. At vero eos et accusam et justo duo dolores et ea rebum. Stet clita kasd gub rgren, no sea takimata sanctus est Lorem ipsum dolor sit amet, consetetur sadipscing elitr, sed diam nonumy eirmod tempor invidunt ut lab ore sit et dolore magna.

\begin{landscape}
\begin{table}[H]
  \caption{2. Satıra geçen örnek tablo adı, 2. Satıra geçen örnek tablo adı, 2. Satıra geçen örnek tablo adı, 2. Satıra geçen örnek tablo adı, 2. Satıra geçen örnek tablo adı.}
  \centering
\begin{tabular}{L{72.6bp}L{76.05bp}L{92.95bp}L{68.65bp}L{68.65bp}L{68.65bp}L{68.65bp}L{68.4bp}}
\hline\hline
\multirow{2}{*}{Parametre} & \multirow{2}{*}{Kolon 2} & \multirow{2}{*}{Kolon 3} & \multicolumn{3}{l|}{\Ts Kolon 4} & \multicolumn{2}{l}{Kolon 5} \\
\cline{4-8}
 &  &  & \Ts Alt kolon & Alt kolon & \multicolumn{1}{l|}{Alt kolon} & Alt kolon & Alt kolon \\
\hline
\Ts Satır 1 & -7.680442 & 7.6986348 & 0.00 & 0.00 & 0.00 & 12 & 12 \\
\Ts Satır 2 & 140 & - & 0.50 & 0.00 & 0.00 & 0 & 0 \\
\Ts Satır 3 & 37.174357 & 37.16192697 & 0.00 & 0.00 & 0.00 & 0 & 24 \\
\Ts Satır 4 & 140 & - & 0.50 & 0.00 & 0.00 & 0 & 0 \\
\Ts Satır 5 & 37.174357 & 37.16192697 & 0.00 & 0.00 & 0.00 & 0 & 24 \\
\Ts Satır 6 & 140 & - & 0.50 & 0.00 & 0.00 & 0 & 0 \\
\Ts Satır 7 & 37.174357 & 37.16192697 & 0.00 & 0.00 & 0.00 & 0 & 24 \\
\Ts Satır 8 & 140 & - & 0.50 & 0.00 & 0.00 & 0 & 0 \\
\Ts Satır 9 & 37.174357 & 37.16192697 & 0.00 & 0.00 & 0.00 & 0 & 24 \\
\Ts Satır 10 & 140 & - & 0.50 & 0.00 & 0.00 & 0 & 0 \\
\Ts Satır 11 & 37.174357 & 37.16192697 & 0.00 & 0.00 & 0.00 & 0 & 24 \\
\Ts Satır 12 & 140 & - & 0.50 & 0.00 & 0.00 & 0 & 0 \\
\Ts Satır 13 & 37.174357 & 37.16192697 & 0.00 & 0.00 & 0.00 & 0 & 24 \\
\Ts Satır 14 & 140 & - & 0.50 & 0.00 & 0.00 & 0 & 0 \\
\Ts Satır 15 & 37.174357 & 37.16192697 & 0.00 & 0.00 & 0.00 & 0 & 24 \\
\hline
\end{tabular}
\end{table}
\begin{table}[H]
  \ContinuedFloat
  \caption[]{\devami{} 2. Satıra geçen örnek tablo adı, 2. Satıra geçen örnek tablo adı, 2. Satıra geçen örnek tablo adı, 2. Satıra geçen örnek tablo adı, 2. Satıra geçen örnek tablo adı.}
  \centering
\begin{tabular}{L{69.0bp}L{72.2bp}L{88.35bp}L{65.15bp}L{65.15bp}L{65.15bp}L{65.15bp}L{64.9bp}}
\hline\hline
\multirow{2}{*}{Parametre} & \multirow{2}{*}{Kolon 2} & \multirow{2}{*}{Kolon 3} & \multicolumn{3}{l|}{\Ts Kolon 4} & \multicolumn{2}{l}{Kolon 5} \\
\cline{4-8}
 &  &  & \Ts Alt kolon & Alt kolon & \multicolumn{1}{l|}{Alt kolon} & Alt kolon & Alt kolon \\
\hline
\Tsd Satır 16 & -7.680442 & 7.6986348 & 0.00 & 0.00 & 0.00 & 12 & 12 \\
\Tsd Satır 17 & 140 & - & 0.50 & 0.00 & 0.00 & 0 & 0 \\
\Tsd Satır 18 & 37.174357 & 37.16192697 & 0.00 & 0.00 & 0.00 & 0 & 24 \\
\Tsd Satır 19 & 140 & - & 0.50 & 0.00 & 0.00 & 0 & 0 \\
\Tsd Satır 20 & 37.174357 & 37.16192697 & 0.00 & 0.00 & 0.00 & 0 & 24 \\
\Tsd Satır 21 & 140 & - & 0.50 & 0.00 & 0.00 & 0 & 0 \\
\Tsd Satır 22 & 37.174357 & 37.16192697 & 0.00 & 0.00 & 0.00 & 0 & 24 \\
\Tsd Satır 23 & 140 & - & 0.50 & 0.00 & 0.00 & 0 & 0 \\
\hline
\end{tabular}
\end{table}

\end{landscape}


\chapter{METİNLER}
\section{Gövde Metinleri}
Lorem ipsum dolor sit amet, consetetur sadipscing elitr, sed diam nonumy eirmod tempor invidunt ut labore et dolore magna aliquyam erat, sed diam voluptua. At vero eos et accusam et justo duo dolores et ea rebum. Stet clita kasd gub rgren, no sea takimata sanctus est Lorem ipsum dolor sit amet, consetetur sadipscing elitr, sed diam nonumy eirmod tempor invidunt ut lab ore sit et dolore magna.

Lorem ipsum dolor sit amet, consetetur sadipscing elitr, sed diam nonumy eirmod tempor invidunt ut labore et dolore magna aliquyam erat, sed diam voluptua. At vero eos et accusam et justo duo dolores et ea rebum. Stet clita kasd gub rgren, no sea takimata sanctus est Lorem ipsum dolor sit amet, consetetur sadipscing elitr, sed diam nonumy eirmod tempor invidunt ut lab ore sit et dolore magna.

\section{Sayfa Kenar Boşlukları}
% [Normal] sp=after:240
Lorem ipsum dolor sit amet, consetetur sadipscing elitr, sed diam Lorem ipsum dolor sit amet, consetetur sadipscing elitr, sed diam nonumy eirmod tempor invidunt ut labore et dolore magna aliquyam erat, sed diam voluptua. At vero eos et accusam et justo duo dolores et ea rebum (Şekil 3.1).

\begin{figure}[H]
  \centering
  \includegraphics[width=12.10cm]{sekiller/image3.png}
  \caption{Sinir hücresi, Çetin (2003)’ten uyarlanmıştır.}
\end{figure}
Sayfa alt boşluğu 2.5 cm’dir. Bundan daha fazla gereksiz boşluk kalmaması gerekir. Sayfalardaki metin, tablo, şekil, vs. bu gözetilerek dengelenmelidir.

Şekiller, tablolar büyütülebilir,  küçültülebilir.

Şekil ve tablolara ait açıklama metinleri (ilk atıf olandan hariç) duruma göre şekil ve tablo öncesine veya sonrasına konulabilir.

Lorem ipsum dolor sit amet, consetetur sadipscing elitr, sed diam nonumy eirmod tempor invidunt ut labore et dolore magna aliquyam erat, sed diam voluptua. At vero eos et accusam et justo duo dolores et ea rebum. Stet clita kasd gub rgren, no sea takimata sanctus est Lorem ipsum dolor sit amet, consetetur sadipscing elitr, sed diam nonumy eirmod tempor invidunt ut lab ore sit et dolore magna.

At vero eos et accusam et justo duo dolores et ea rebum. Stet clita kasd gub rgren, no sea takimata sanctus est Lorem ipsum dolor sit amet, consetetur sadipscing elitr, sed diam nonumy eirmod tempor invidunt ut lab ore sit et dolore magna.

Lorem ipsum dolor sit amet, consetetur sadipscing elitr, sed diam nonumy eirmod tempor invidunt ut labore et dolore magna aliquyam erat, sed diam voluptua. At vero eos et accusam et justo duo dolores et ea rebum. Stet clita kasd gub rgren, no sea takimata sanctus est Lorem ipsum dolor sit amet, consetetur sadipscing elitr, sed diam nonumy eirmod tempor invidunt ut lab ore sit et dolore magna.

\section{Denklemler}
Lorem ipsum dolor sit amet, consetetur sadipscing elitr, sed diam nonumy eirmod tempor invidunt ut labore et dolore magna aliquyam erat, sed diam voluptua. At vero eos et accusam et justo duo dolores et ea rebum. Stet clita kasd gub rgren, no sea takimata sanctus est Lorem ipsum dolor sit amet, consetetur sadipscing elitr, sed diam nonumy eirmod tempor invidunt ut lab ore sit et dolore magna (3.1).

\begin{denklem}
  \symup{φ=ω-{β}_{k+1}}
  \label{denk:3.1}
\end{denklem}
\begin{denklem}
  \symup{\lim\nolimits_{n→∞}S({x}_{n},{x}_{n},{x}_{n+1})=0}
  \label{denk:3.2}
\end{denklem}
Parametreler tek tek açıklanır. Denklem 3.1’de, 3.1’de veya formül 3.1’de görüleceği üzere. Lorem ipsum dolor sit amet, consetetur sadipscing elitr, sed diam nonumy eirmod tempor invidunt ut labore et dolore denklem 3.1’in magna aliquyam erat.

Denklemler denklem editörüyle ve aşağıdaki kurallara uygun olarak yazılır.

Denklemler 1,5 satır aralıkla yazılır.

Denklemlerden önce ve sonra 6 punto (nk) aralık bırakılır.

Denklemlerden önce ve sonra boş satır bırakılmamalıdır. Denklemler metin bloğuna ortalı olarak hizalanır.

Denklemlere, ilgili bölüm içinde bölüm numarası ilk numara olacak şekilde, 1’den başlayarak sıra ile numara verilir. Bu numaralar normal yazı karakteri ile yazılmalı (Times New Roman, 12 Punto), koyu (bold) veya italik olmamalıdır. Örnek olarak birinci bölümde (1.1), (1.2), ikinci bölümde (2.1), (2.2) veya gerekli olması halinde aynı denklemin alt ifadeleri (1.1a) ve (1.1b) şeklinde denklemin bulunduğu satırın en sağına yazılır. Denklem numarasının sonuna nokta konulmamalıdır. Ekler bölümünde verilen denklemler, verildikleri bölüm belirtilerek numaralandırılır. (Örnek: A.1.1, A.1.2).

\subsection{Üçüncü derece başlık örneği}
% [GOVDE] sp=after:240
Lorem ipsum dolor sit amet, consetetur sadipscing elitr, sed diam nonumy eirmod tempor invidunt ut labore et dolore magna aliquyam erat, sed diam voluptua. At vero eos et accusam et justo duo dolores et ea rebum. Stet clita kasd gub rgren, no sea takimata sanctus est Lorem ipsum dolor sit amet, consetetur sadipscing elitr, sed diam nonumy eirmod tempor invidunt ut lab ore sit et dolore magna.

\begin{figure}[H]
  \centering
  \includegraphics[width=8.11cm]{sekiller/image4.png}
  \caption{Birden fazla satırlı şekil isimlendirmesinde örnek, birden fazla satırlı şekil isimlendirmesinde örnek.}
\end{figure}
Lorem ipsum dolor sit amet, consetetur sadipscing elitr, sed diam nonumy eirmod tempor invidunt ut labore et dolore magna aliquyam erat, sed diam voluptua. At vero eos et accusam et justo duo dolores et ea rebum. Stet clita kasd gub rgren, no sea takimata sanctus est Lorem ipsum dolor sit amet, consetetur sadipscing elitr, sed diam nonumy eirmod tempor invidunt ut lab ore sit et dolore magna.

Lorem ipsum dolor sit amet, consetetur sadipscing elitr, sed diam nonumy eirmod tempor invidunt ut labore et dolore magna aliquyam erat, sed diam voluptua. At vero eos et accusam et justo duo dolores et ea rebum. Stet clita kasd gub rgren, no sea takimata sanctus est Lorem ipsum dolor sit amet, consetetur sadipscing elitr, sed diam nonumy eirmod tempor invidunt ut lab ore sit et dolore magna.

\subsubsection{Dördüncü derece başlık örneği dördüncü derece başlık örneği dördüncü derece başlık örneği}
% [GOVDE] sp=after:240
Lorem ipsum dolor sit amet, consetetur sadipscing elitr, sed diam nonumy eirmod tempor invidunt ut labore et dolore magna aliquyam erat, sed diam voluptua. At vero eos et accusam et justo duo dolores et ea rebum. Stet clita kasd gub rgren, no sea takimata sanctus est Lorem ipsum dolor sit amet, consetetur sadipscing elitr, sed diam nonumy eirmod tempor invidunt ut lab ore sit et dolore magna (3.2). Lorem ipsum dolor sit amet, consetetur sadipscing elitr, sed diam nonumy eirmod tempor invidunt ut labore et dolore magna aliquyam erat, sed diam voluptua. At vero eos et accusam et justo duo dolores et ea rebum. Lorem ipsum dolor sit amet, consetetur sadipscing elitr, sed diam nonumy eirmod tempor invidunt ut labore et dolore magna aliquyam erat, sed diam voluptua. At vero eos et accusam et justo duo dolores et ea rebum. Lorem ipsum dolor sit amet, consetetur sadipscing elitr, sed diam nonumy eirmod tempor invidunt ut labore et dolore magna aliquyam erat, sed diam voluptua. At vero eos et accusam et justo duo dolores et ea rebum.

\begin{figure}[H]
  \centering
  \includegraphics[width=11.23cm]{sekiller/image5.png}
  \caption{Örnek şekil ismi nokta ile bitirilmelidir.}
\end{figure}
\begin{denklem}
  \symup{{u}^{∆}=g(t,u)}
  \label{denk:3.3}
\end{denklem}
Lorem ipsum dolor sit amet, consetetur sadipscing elitr, sed diam nonumy eirmod tempor invidunt ut labore et dolore magna aliquyam erat, sed diam voluptua. Lorem ipsum dolor sit amet, consetetur sadipscing elitr, sed diam nonumy eirmod tempor invidunt ut labore et dolore magna aliquyam erat, sed diam voluptua. At vero eos et accusam et justo duo dolores et ea rebum.

\section{İleri Arıtma Metotları}
Lorem ipsum dolor sit amet, consetetur sadipscing elitr, sed diam nonumy eirmod tempor invidunt ut labore et dolore magna aliquyam erat, sed diam voluptua. At vero eos et accusam et justo duo dolores et ea rebum. Stet clita kasd gub rgren, no sea takimata sanctus est Lorem ipsum dolor sit amet, consetetur sadipscing elitr, sed diam nonumy eirmod tempor invidunt ut lab ore sit et dolore magna.

Lorem ipsum dolor sit amet, consetetur sadipscing elitr, sed diam nonumy eirmod tempor invidunt ut labore et dolore magna aliquyam erat, sed diam voluptua. At vero eos et accusam et justo duo dolores et ea rebum. Stet clita kasd gub rgren, no sea takimata sanctus est Lorem ipsum dolor sit amet, consetetur sadipscing elitr, sed diam nonumy eirmod tempor invidunt ut lab ore sit et dolore magna.

Lorem ipsum dolor sit amet, consetetur sadipscing elitr, sed diam nonumy eirmod tempor invidunt ut labore et dolore magna aliquyam erat, sed diam voluptua. At vero eos et accusam et justo duo dolores et ea rebum. Stet clita kasd gub rgren, no sea takimata sanctus est Lorem ipsum dolor sit amet, consetetur sadipscing elitr, sed diam nonumy eirmod tempor invidunt ut lab ore sit et dolore magna.

Lorem ipsum dolor sit amet, consetetur sadipscing elitr, sed diam nonumy eirmod tempor invidunt ut labore et dolore magna aliquyam erat, sed diam voluptua. At vero eos et accusam et justo duo dolores et ea rebum. Stet clita kasd gub rgren, no sea takimata sanctus est Lorem ipsum dolor sit amet, consetetur sadipscing elitr, sed diam nonumy eirmod tempor invidunt ut lab ore sit et dolore magna.

\section{Deneysel Çalışmalar}
Lorem ipsum dolor sit amet, consetetur sadipscing elitr, sed diam nonumy eirmod tempor invidunt ut labore et dolore magna aliquyam erat, sed diam voluptua. At vero eos et accusam et justo duo dolores et ea rebum. Stet clita kasd gub rgren, no sea takimata sanctus est Lorem ipsum dolor sit amet, consetetur sadipscing elitr, sed diam nonumy eirmod tempor invidunt ut lab ore sit et dolore magna (Nelson, 1988).

Lorem ipsum dolor sit amet, consetetur sadipscing elitr, sed diam nonumy eirmod tempor invidunt ut labore et dolore magna aliquyam erat, sed diam voluptua. At vero eos et accusam et justo duo dolores et ea rebum. Lorem ipsum dolor sit amet, consetetur sadipscing elitr, sed diam nonumy eirmod tempor invidunt ut labore et dolore magna aliquyam erat, sed diam voluptua.

Lorem ipsum dolor sit amet, consetetur sadipscing elitr, sed diam nonumy eirmod tempor invidunt ut labore et dolore magna aliquyam erat, sed diam voluptua. At vero eos et accusam et justo duo dolores et ea rebum. Stet clita kasd gub rgren, no sea takimata sanctus est Lorem ipsum dolor sit amet, consetetur sadipscing elitr, sed diam nonumy eirmod tempor invidunt ut lab ore sit et dolore magna (Nelson, 1988).

Lorem ipsum dolor sit amet, consetetur sadipscing elitr, sed diam nonumy eirmod tempor invidunt ut labore et dolore magna aliquyam erat, sed diam voluptua. At vero eos et accusam et justo duo dolores et ea rebum. Lorem ipsum dolor sit amet, consetetur sadipscing elitr, sed diam nonumy eirmod tempor invidunt ut labore et dolore magna aliquyam erat, sed diam voluptua.


\chapter{ATIFLAR, ALINTILAR VE DİPNOTLAR}
Bu bölümde atıflar, alıntılar ve dipnotların nasıl olması gerektiği hakkında bilgi verilecektir.

\section{Atıflar (kaynakların metin içinde gösterimi)}
\subsection{Yazar soyadına göre atıf verme}
Kaynaklar metin içinde yazar soyadı ve tarih belirtilerek verilir. Kaynaklar sayfasında yazar soyadına göre alfabetik olarak sıralama yapılır. Metin içinde kaynak, cümlenin başlangıcında veya içinde verilecekse, Baran (2003) şeklinde, kaynak cümle sonunda verilecekse (Baran, 2003). şeklinde gösterilir. Nokta işareti kaynaktan hemen sonra konulur.

Kaynak birden fazla yazara ait olduğunda, yazar sayısı iki ise, cümle başında veya içinde Uyar ve Yılmaz (2020) şeklinde, cümle sonunda ise (Uyar ve Yılmaz, 2020). şeklinde yazılır. Yazar sayısı ikiden fazla ise cümle başında veya içinde Uyar ve ark. (2018) şeklinde, cümle sonunda ise (Uyar ve ark, 2018). şeklinde yazılır.

Aynı yazara ait ve aynı yıl içinde yayınlanmış yayınlar Uyar (2008a), Uyar (2008b) şeklinde numaralandırılır. Aynı parantez içerisinde aynı yazarın 2 ve daha fazla eserine atıfta bulunma; son yayınlanan eseri en son belirterek aynı parantez içerisinde gösterilebilirler. Örneğin; Son araştırmalar \ldots{}. göstermiştir (Uyar, 2008, 2020).

Aynı parantez içerisinde 2 ya da daha fazla esere atıf, yazar soyadlarına göre alfabetik sıraya uygun şekilde verilir; (Ateş, 2021; Çelik, 2012; Yılmaz, 2022).

İkincil kaynak drumunda (atıf yapılan kaynak başka bir kaynağa atıfta bulunuyorsa), metinde orijinal kaynağa atıfta bulunulur ve parantez içerisinde orijinal kaynağa atıfta bulunan yazara gönderme yapılır. Referans listesinde sadece orijinal kaynağa atıfta bulunan kaynak için giriş yapılır; orijinal kaynak için referans girişi yapılmaz. Kaynaklar bölümünde künyeler yazar soyadına göre sıralanır.

\subsection{Numara ile atıf verme}
Metin içinde [ ] köşeli parantez içinde numaralandırılır. Tezde ilk verilen kaynak [1] numara ile başlar ve veriliş sırasına göre numaralandırılır. Kaynaklara metin içerisinde aşağıdaki biçimlerde atıf yapılır.

[1]\quad{}\quad{}1 nolu kaynak,

[3-7]\quad{}\quad{}3 ve 7 arası (3, 4, 5, 6 ve 7 nolu ) kaynaklar,

[3, 6]\quad{}\quad{}3 ve 6 nolu kaynaklar,

[3, 5, 8]\quad{}3, 5 ve 8 nolu kaynaklar,

[3, 6-9]\quad{}3 ve 6 ile 9 nolu kaynaklar arasındaki kaynaklar.

Aynı isimli birden fazla cildi olan kaynakların, kullanılan her bir cildine ayrı kaynak numarası verilmelidir.

\subsubsection{Kaynakların listelenmesi}
Kaynaklar bölüm başlığı 1 satır aralığı kullanılarak diğer birinci başlıklardaki gibi (önce 72, sonra 18 punto) aralık bırakılarak yazılmalıdır. Kaynaklar ise tek satır aralıklı, önce ve sonra 6 punto (nk) boşluk bırakılarak yazılır.

Tezde atıf yapılan bütün kaynaklar ‘Kaynaklar’ bölümünde listelenmelidir.

Yazar soyadına göre atıf yapılan tezlerde kaynaklar listesi alfabetik olarak sıralanır. Numaralandırma yöntemine göre atıf yapılan tezlerde ise kaynaklar tezde veriliş sırasına göre numaralanır ve sıralanır.

Kaynaklar bölümünde aynı yazarın birden çok eseri yer alıyorsa, eserler, yayın yılına göre eskiden yeniye doğru sıralanır. Aynı yazarın iki farklı eserinin yayın tarihleri aynıysa, kaynaklar listesindeki sıralama, künyede bir sonraki öğe olan eser adına göre yapılır. Atıflarda ayrımı sağlamak için tarih bilgisinin yanına a’dan başlayan harfler eklenir.

Aynı yazarın tek yazarlı eserleri çok yazarlı eserlerinden önce sıralanır. Çok yazarlı iki eserin ilk yazarları aynıysa, ikinci yazarın soyadı, ikinci yazarlar da aynıysa üçüncü yazarın soyadı alfabetik düzende belirleyicidir.

Aynı soyadını taşıyan iki farklı yazarın eserleri adlarına göre alfabetik sıraya girer.

Yazar bir tüzel kuruluşsa, eser, yazar konumundaki tüzel kuruluşun adıyla alfabetik listeye girer.

Bir eserin yazarı veya editörü yoksa künye eser adına hazırlanacağı için, kaynak, eser adından alfabetik listeye girer. Eser adı rakamla başlıyorsa sıralamada rakamın okunuşu dikkate alınır.

İkincil kaynak durumunda (atıf yapılan kaynak başka bir kaynağa atıfta bulunuyorsa), metinde orijinal kaynağa atıfta bulunulur ve parantez içerisinde orijinal kaynağa atıfta bulunan yazara gönderme yapılır. Referans listesinde sadece orijinal kaynağa atıfta bulunan kaynak için giriş yapılır; orijinal kaynak için referans girişi yapılmaz. Kaynaklar bölümünde künyeler yazar soyadına göre sıralanır.

İkincil kaynak, ilgili kaynağa ulaşılamadığı durumlarda “aktaran” ifadesi ile yapılmalıdır. Örneğin;

Parantez içerisinde kullanımı:

(Baran, 1945; akt. Çoruhlu, 2015).

Metin içi kullanımı:

Baran (aktaran Çoruhlu, 2015) 1945’ te \ldots{}

\section{Alıntılar}
Genel olarak alıntılar kelime, imla ve noktalama bakımından aslına uygun olarak yapılır. Alıntı yapılan parçada bir yanlış varsa, doğrusu köşeli parantez içerisinde belirtilmek koşuluyla metin aynen nakledilir.

Kırk kelimeden daha az uzunluktaki kısa alıntılar çift tırnak içerisinde verilir. Alıntının sonunda ilgili kaynağa atıf yapılıp atıftan sonra nokta koyulur.

Kırk kelimeden fazla olan uzun alıntılar tırnak içerisinde gösterilmezler. Uzun alıntılar soldan 1 sekme (1,27 cm) içerden verilir. İçerden verilen uzun alıntılarda, 2 yazı karakteri daha küçük karakter kullanılır. Eğer alıntı yapılan metinde aynı cümle içinde kelime atlanarak yazılıyorsa bu üç nokta (...) ile, eğer bazı cümleler atlanıyorsa bu dört nokta (....) ile belirtilmelidir. Ancak, çok sık ve çok uzun alıntılardan kaçınılması tavsiye edilir. Kısa alıntılardan farklı olarak noktalama atıftan sonra değil de önce yapılır. Örneğin; . (Yılmaz, 2021, s. 196)

40 kelimeden fazla olan alıntı örneği;

Ana metin ana metin ana metin ana metin ana metin ana metin ana metin ana metin ana metin ana metin. Bu görüşe karşı olan bazı çalışmalarda bulunmaktadır:

\begin{alinti}
Bilim, araştırma ve geliştirme faaliyetleri ekonomik kalkınma ve sosyal gelişmenin gerçekleşmesinde etkin araçlar olarak belirmektedir. Bilimsel ve teknolojik araştırmaların üretim ve karar verme mekanizmalarına olan etkileri özellikle gelişmiş ülkelerde hemen her alanda kendini belli etmektedir. Oysa az gelişmiş ülkelerde bu ilişkiyi fark etmek bile zordur. (Eryaman, 1973, 685)
\end{alinti}

Betimleme yöntemi, Kaptan (1995) tarafından şu biçimde tanımlanmıştır:

\begin{alinti}
Olayların, objelerin, varlıkların, kurumların, grupların ve çeşitli alanların ne olduğunu betimlemeye, açıklamaya çalışan incelemelerdir. Bunlar nedir? sorusuna cevap bulmaya yöneliktir. Bununla mevcut durumlar, koşullar, özellikler aynen ortaya konmaya çalışılır. Betimleme araştırmaları, mevcut olayların daha önceki olay ve koşullarla ilişkilerini de dikkate alarak, durumlar arasındaki etkileşimi açıklamayı hedefler (s. 59).
\end{alinti}

Bazıları bu fikre karşı çıktılar:

\begin{alinti}
Aynı anda, aynı yerde bulunmak bütün grup üyeleri arasında samimi bir etkileşim olmasını gerektirmez. Bir ritüeli gerçekleştirmek ya da bir olayı kutlamak için aynı yerde, bir araya toplanan yüzlerce, hatta binlerce kişiyi düşünün. Bu ortamlarda, katılımcılar grubun görsel tezahürünü yaşayıp, fiziksel olarak bir arada olsa da çevrelerindeki kişilerle doğrudan, samimi bir bağlantı kurma şansları grubun büyüklüğüne bağlıdır (Purcell, 1997, ss. 111-112).
\end{alinti}

40 kelimeden uzun olan, kelime ve cümle atlanarak kullanılarak yazılan alıntı örneği:

\begin{alinti}
Bilim, dünyamızda olup biten olguları betimleme ve açıklama yoluyla anlama girişimidir. Olguları betimleme, onları saptama, sınıflama ve dile getirme gibi işlemleri kapsar. .... Bilimin açıklama yönüne gelince bu konu bizi çok daha geniş sorunlara götürecek niteliktedir. ... Hipotez, doğa yasası, teori, nedensellik ve olasılık ilkeleri gibi kavramları ele almaya ihtiyaç vardır (Yıldırım, 1991, 95).
\end{alinti}

Cümle başındaki alıntı örnekleri;

Orhan (2018)’e göre, “Bu tür çalışmalar çok sayıda fiziksel, kimyasal ve biyolojik parametrenin çok sayıda ölçüm noktasında incelenmesini gerektirmektedir” (s. 199).

Ertürk (1998) eğitimi, “bireyin davranışlarında kendi yaşantısı yoluyla ve kasıtlı olarak istendik değişiklik meydana getirme süreci” olarak tanımlamaktadır (Ertürk, 1998, 12).Bilimi çeşitli şekillerde tanımlamak mümkündür. Örneğin, Yıldırım (1991) bilimi “dünyamızda olup biten olguları ‘betimleme ve açıklama’ yoluyla anlama girişimidir” şeklinde tanımlıyor (s. 95).

Cümle arasındaki kısa alıntı örneği;

Bu sonuçları yorumlarken, Robbins ve ark. (2003) “terapistler, bırakılma durumda istemeden de olsa, ergenin ihtiyaçlarını ya da kaygılarını tam olarak karşılamayabilirler” (s. 541) diyerek katkıda bulundular.

Cümle sonundaki kısa alıntı örneği;

Bu konuda kafa karışıklığına sebep olan şey, palyatif bakım sırasındaki rollerin doğasıdır, “tıbbi gereksinimler tıp alanındaki görevlilerce karşılanır; tıbbi olmayan ihtiyaçlar ise ekipteki herhangi birisi tarafından karşılanabilir” (Csikai ve Chaitin, 2006, s. 112).

\section{Dipnotlar}
Tezlerde içeriği genişletici, güçlendirici veya ilave nitelikteki bilgiler (içerik dipnotu) kullanılabilir\footnote{Dipnotlar ile kaynak gösterimi yapılmaz. Dipnotlar tez içerisinde içeriği genişletici, güçlendirci veya ilave nitelikteki bilgileri vermek için kullanılır. Verilen genişletici, güçlendirci veya ilave nitelikteki bilgiler zorunlulukla kaynak içeriyorsa bu kaynak mutlaka kaynaklar bölümünde verilmelidir.}.

Tezlerde içeriği genişletici, güçlendirici veya ilave nitelikteki bilgiler (içerik dipnotu) kullanılabilir.

Dipnot numaraları alıntının hemen sonuna koyulur. Alıntı paragrafsa dipnot numarası paragrafın son kelimesinin üzerine, alıntı bir kavram veya isimse, kavram veya ismin hemen üzerine yazılır.

Metin içerisindeki dipnot numarası; satır hizasının üzerinde1 şeklinde görünür olmalıdır. Numara sonrasında herhangi bir noktalama işareti konmamalıdır.

Dipnot, ilgili sayfanın altına metinden 2 karakter küçük (10 punto) yazı ile yazılmalıdır.

Dipnot çizgisi ile dipnot numarası arasında bir aralık; dipnot numarası ile dipnotun ilk satırı arasında ise yarım aralık bırakılmalıdır. Dipnotlar metinden ince yatay bir çizgi ile ayrılmalıdır.

\section{İkinci Derece Başlık Nasıl: İlk Harfler Büyük}
Lorem ipsum dolor sit amet, consetetur sadipscing elitr, sed diam nonumy eirmod tempor invidunt ut labore et dolore magna aliquyam erat, sed diam voluptua. At vero eos et accusam et justo duo dolores et ea rebum. Stet clita kasd gub rgren, no sea

\subsection{Üçüncü derece başlık nasıl: ilk harf büyük diğerleri küçük}
Lorem ipsum dolor sit amet, consetetur sadipscing elitr, sed diam nonumy eirmod tempor invidunt ut labore et dolore magna aliquyam erat, sed diam voluptua. At vero eos et accusam et justo duo dolores et ea rebum. Stet clita kasd gub rgren, no sea

\subsubsection{Dördüncü derece başlık nasıl: ilk harf büyük diğerleri küçük}
Stet clita kasd gub rgren, no sea takimata sanctus est Lorem ipsum dolor sit amet, consetetur sadipscing elitr, sed diam nonumy eirmod tempor invidunt ut lab ore sit et dolore magna.

\baslikbes{Beşinci derece başlık: dördüncü dereceden sonrası numaralandırılmaz}
% [Normal] sp=after:240
Stet clita kasd gub rgren, no sea takimata sanctus est Lorem ipsum dolor sit amet, consetetur sadipscing elitr, sed diam nonumy eirmod tempor invidunt ut lab ore sit et dolore magna.

\begin{figure}[H]
  \centering
  \includegraphics[width=9.33cm]{sekiller/image6.png}
  \caption{Örnek şekil.}
\end{figure}
Stet clita kasd gub rgren, no sea takimata sanctus est Lorem ipsum dolor sit amet, consetetur sadipscing elitr, sed diam nonumy eirmod tempor invidunt ut lab ore sit et dolore magna.

\begin{table}[H]
  \caption{Tablo örneği.}
  \centering
\begin{tabular}{C{3.05cm}C{2.10cm}C{2.16cm}C{2.52cm}}
\hline\hline
\Ns Kolon A & Kolon B & Kolon C & Kolon D \\
\hline
\Ns Satır A & Satır A & Satır A & Satır A \\
\Ns Satır B & Satır B & Satır B & Satır B \\
\Ns Satır C & Satır C & Satır C & Satır C \\
\hline
\end{tabular}
\end{table}
% [Tablodansonragelenmetin] 
Stet clita kasd gub rgren, no sea takimata sanctus est Lorem ipsum dolor sit amet, consetetur sadipscing elitr, sed diam nonumy eirmod tempor invidunt ut lab ore sit et dolore magna. Stet clita kasd gub rgren, no sea takimata sanctus est Lorem ipsum dolor sit amet, consetetur sadipscing elitr, sed diam nonumy eirmod tempor invidunt ut lab ore sit et dolore magna.

Stet clita kasd gub rgren, no sea takimata sanctus est Lorem ipsum dolor sit amet, consetetur sadipscing elitr, sed diam nonumy eirmod tempor invidunt ut lab ore sit et dolore magna.


\chapter{GEREKLİ İSE BÖLÜM BAŞLIĞI}
Lorem ipsum dolor sit amet, consetetur sadipscing elitr, sed diam nonumy eirmod tempor invidunt ut labore et dolore magna aliquyam erat, sed diam voluptua. At vero eos et accusam et justo duo dolores et ea rebum. Stet clita kasd gub rgren, no sea takimata sanctus est Lorem ipsum dolor sit amet, consetetur sadipscing elitr, sed diam nonumy eirmod tempor invidunt ut lab ore sit et dolore magna.

\section{Deney Sonuçlar}
In this thesis, the necessary steps for constructing an end-to-end streamflow forecasting system were discussed. These steps include the use

\section{İkinci Derece Başlık Nasıl: İlk Harfler Büyük}
Lorem ipsum dolor sit amet, consetetur sadipscing elitr, sed diam nonumy eirmod tempor invidunt ut labore et dolore magna aliquyam erat, sed diam voluptua. At vero eos et accusam et justo duo dolores et ea rebum. Stet clita kasd gub rgren, no sea

\subsection{Üçüncü derece başlık nasıl: ilk harf büyük diğerleri küçük}
Lorem ipsum dolor sit amet, consetetur sadipscing elitr, sed diam nonumy eirmod tempor invidunt ut labore et dolore magna aliquyam erat, sed diam voluptua. At vero eos et accusam et justo duo dolores et ea rebum. Stet clita kasd gub rgren, no sea

\subsubsection{Dördüncü derece başlık nasıl: ilk harf büyük diğerleri küçük}
Stet clita kasd gub rgren, no sea takimata sanctus est Lorem ipsum dolor sit amet, consetetur sadipscing elitr, sed diam nonumy eirmod tempor invidunt ut lab ore sit et dolore magna.

\baslikbes{Beşinci derece başlık nasıl: ilk harf büyük diğerleri küçük}
% [GOVDE] sp=after:240
Lorem ipsum dolor sit amet, consetetur sadipscing elitr, sed diam nonumy eirmod tempor invidunt ut labore et dolore magna aliquyam erat, sed diam voluptua.

\begin{figure}[H]
  \centering
  \includegraphics[width=11.11cm]{sekiller/image1.png}
  \caption{Beşinci bölümde örnek şekil.}
\end{figure}
Stet clita kasd gub rgren, no sea takimata sanctus est Lorem ipsum dolor sit amet, consetetur sadipscing elitr, sed diam nonumy eirmod tempor invidunt ut lab ore sit et dolore magna.

\begin{table}[H]
  \caption{Beşinci bölümde örnek tablo.}
  \centering
\begin{tabular}{C{3.05cm}C{2.10cm}C{2.16cm}C{2.52cm}}
\hline\hline
\Ns Kolon A & Kolon B & Kolon C & Kolon D \\
\hline
\Ns Satır A & Satır A & Satır A & Satır A \\
\Ns Satır B & Satır B & Satır B & Satır B \\
\Ns Satır C & Satır C & Satır C & Satır C \\
\hline
\end{tabular}
\end{table}
% [Normal] sp=before:240
Stet clita kasd gub rgren, no sea takimata sanctus est Lorem ipsum dolor sit amet, consetetur sadipscing elitr, sed diam nonumy eirmod tempor invidunt ut lab ore sit et dolore magna. Stet clita kasd gub rgren, no sea takimata sanctus est Lorem ipsum dolor sit amet, consetetur sadipscing elitr, sed diam nonumy eirmod tempor invidunt ut lab ore sit et dolore magna.

Stet clita kasd gub rgren, no sea takimata sanctus est Lorem ipsum dolor sit amet, consetetur sadipscing elitr, sed diam nonumy eirmod tempor invidunt ut lab ore sit et dolore magna. Stet clita kasd gub rgren, no sea takimata sanctus est Lorem ipsum dolor sit amet, consetetur sadipscing elitr, sed diam nonumy eirmod tempor invidunt ut lab ore sit et dolore magna. 
\clearpage


\chapter{SONUÇ VE ÖNERİLER}
Lorem ipsum dolor sit amet, consetetur sadipscing elitr, sed diam nonumy eirmod tempor invidunt ut labore et dolore magna aliquyam erat, sed diam voluptua. At vero eos et accusam et justo duo dolores et ea rebum. Stet clita kasd gub rgren, no sea takimata sanctus est Lorem ipsum dolor sit amet, consetetur sadipscing elitr, sed diam nonumy eirmod tempor invidunt ut lab ore sit et dolore magna.

\section{Teorik Modeller}
Lorem ipsum dolor sit amet, consetetur sadipscing elitr, sed diam nonumy eirmod tempor invidunt ut labore et dolore magna aliquyam erat, sed diam voluptua. At vero eos et accusam et justo duo dolores et ea rebum. Stet clita kasd gub rgren, no sea takimata sanctus est Lorem ipsum dolor sit amet, consetetur sadipscing elitr, sed diam nonumy eirmod tempor invidunt ut lab ore sit et dolore magna.

\section{İkinci Derece Başlık Nasıl: İlk Harfler Büyük}
Lorem ipsum dolor sit amet, consetetur sadipscing elitr, sed diam nonumy eirmod tempor invidunt ut labore et dolore magna aliquyam erat, sed diam voluptua. At vero eos et accusam et justo duo dolores et ea rebum. Stet clita kasd gub rgren, no sea

\subsection{Üçüncü derece başlık nasıl: ilk harf büyük diğerleri küçük}
Lorem ipsum dolor sit amet, consetetur sadipscing elitr, sed diam nonumy eirmod tempor invidunt ut labore et dolore magna aliquyam erat, sed diam voluptua. At vero eos et accusam et justo duo dolores et ea rebum. Stet clita kasd gub rgren, no sea

\subsubsection{Dördüncü derece başlık nasıl: ilk harf büyük diğerleri küçük}
% [GOVDE] sp=after:240
Stet clita kasd gub rgren, no sea takimata sanctus est Lorem ipsum dolor sit amet, consetetur sadipscing elitr, sed diam nonumy eirmod tempor invidunt ut lab ore sit et dolore magna.

\clearpage
\begin{figure}[H]
  \centering
  \includegraphics[width=11.11cm]{sekiller/image1.png}
  \caption{Altıncı bölümde örnek şekil.}
\end{figure}
Lorem ipsum dolor sit amet, consetetur sadipscing elitr, sed diam nonumy eirmod tempor invidunt ut labore et dolore magna aliquyam erat, sed diam voluptua. At vero eos et accusam et justo duo dolores et ea rebum. Stet clita kasd gub rgren, no sea.

\begin{table}[H]
  \caption{Altıncı bölümde bir tablo.}
  \centering
\begin{tabular}{C{3.05cm}C{2.10cm}C{2.16cm}C{2.52cm}}
\hline\hline
\Ns Kolon A & Kolon B & Kolon C & Kolon D \\
\hline
\Ns Satır A & Satır A & Satır A & Satır A \\
\Ns Satır B & Satır B & Satır B & Satır B \\
\Ns Satır C & Satır C & Satır C & Satır C \\
\hline
\end{tabular}
\end{table}
% [GOVDE] sp=before:240
Stet clita kasd gub rgren, no sea takimata sanctus est Lorem ipsum dolor sit amet, consetetur sadipscing elitr, sed diam nonumy eirmod tempor invidunt ut lab ore sit et dolore magna. Stet clita kasd gub rgren, no sea takimata sanctus est Lorem ipsum dolor sit amet, consetetur sadipscing elitr, sed diam nonumy eirmod tempor invidunt ut lab ore sit et dolore magna. Stet clita kasd gub rgren, no sea takimata sanctus est Lorem ipsum dolor sit amet, consetetur sadipscing elitr, sed diam nonumy eirmod tempor invidunt ut lab ore sit et dolore magna.

Stet clita kasd gub rgren, no sea takimata sanctus est Lorem ipsum dolor sit amet, consetetur sadipscing elitr, sed diam nonumy eirmod tempor invidunt ut lab ore sit et dolore magna.


\bolumsuz{KAYNAKLAR}
Kaynaklar metin içinde yazar soyadı ve tarih belirtilerek verilir. Kaynaklar sayfasında yazar soyadına göre alfabetik olarak sıralama yapılır. Metin içinde kaynak, cümlenin başlangıcında veya içinde verilecekse, Baran (2013) şeklinde, kaynak cümle sonunda verilecekse (Baran, 2013). şeklinde gösterilir. Nokta işareti kaynaktan hemen sonra konulur.

Kaynak birden fazla yazara ait olduğunda, yazar sayısıi ise, cümle başında veya içinde Uyar ve Yılmaz (2020) şeklinde, cümle sonunda ise (Uyar ve Yılmaz, 2020). şeklinde yazılır. Yazar sayısı ikiden fazla ise cümle başında veya içinde Uyar ve ark. (2018) şeklinde, cümle sonunda ise (Uyar ve ark, 2018). şeklinde yazılır.

Aynı yazara ait ve aynı yıl içinde yayımlanmış yayınlar Uyar (2008a), Uyar (2008b) şeklinde numaralandırılır. Aynı parantez içerisinde aynı yazarın 2 ve daha fazla eserine atıfta bulunma; son yayınlanan eseri en son belirterek aynı parantez içerisinde gösterilebilirler. Örneğin; Son araştırmalar \ldots{}. göstermiştir (Uyar, 2008, 2020).

Aynı parantez içerisinde 2 ya da daha fazla esere atıf yazar soyadlarına göre alfabetik sıraya uygun şekilde verilir; (Ateş, 2021; Çelik, 2012; Yılmaz, 2022).

İkincil kaynak durumunda (atıf yapılan kaynak başka bir kaynağa atıfta bulunuyorsa), metinde orijinal kaynağa atıfta bulunulur ve parantez içerisinde orijinal kaynağa atıfta bulunan yazara gönderme yapılır. Referans listesinde sadece orijinal kaynağa atıfta bulunan kaynak için giriş yapılır; orijinal kaynak için referans girişi yapılmaz. Kaynaklar bölümünde künyeler yazar soyadına göre sıralanır.

Başka bir kaynaktan alıntı yapma (İkincil Kaynaklar; Aktaran):

İkincil kaynak, ilgili kaynağa ulaşılamadığı durumlarda “aktaran” ifadesi ile yapılmalıdır. Örneğin;

Parantez içerisinde kullanımı:

(Baran, 1945; akt. Çoruhlu, 2015).

Metin içi kullanımı:

Baran (aktaran Çoruhlu, 2015) 1945’ te \ldots{}

\textbf{APA 7  formatına göre kaynakça yazım örnekleri}

\kaynak{\notrenk{0070C0}{\textbf{KİTAP}}}

{\teksatir \textbf{Kitap: Tek yazarlı }\par}

\kaynak{\notrenk{FF0000}{Yazarın soyadı, Yazarın adının baş harfi. (Yıl). }\notrenk{FF0000}{\textit{Kitabın adı }}\notrenk{FF0000}{(Baskı Sayısı). Yayınevi.}}

\kaynak{[Metin içinde atıf: (Jackson, 2019; Sapolsky, 2017; Tarhan, 2012]}

\kaynak{Jackson, L. M. (2019). \textit{The psychology of prejudice: From attitudes to social action} (2nd ed.). American Psychological. Association. https://doi.org/10.1037/0000168-000}

\kaynak{Sapolsky, R. M. (2017). \textit{Behave: The biology of humans at our best and worst}. Penguin Books.}

\kaynak{Tarhan, N. (2012). \textit{Mesnevi terapi}. Timaş Yayınları.}

{\teksatir \textbf{Kitap: İki yazarlı }\par}

{\teksatir [Metin içinde atıf: (Nurmedov ve Tarhan, 2012; Svendsen ve Løber, 2020)]\par}

\kaynak{Nurmedov, S. ve Tarhan, N. (2012). \textit{Bağımlılık: Sanal veya Gerçek }(4. Baskı). Timaş Yayınları.}

\kaynak{Svendsen, S., and Løber, L. (2020). \textit{The big picture/Academic writing: The one-hour guide} (3rd digital ed.). Hans Reitzel Forlag.}

\kaynak{\textbf{Kitap: 3 veya daha fazla yazarlı  }}

\kaynak{[Metin içinde atıf: (Ülkü Arıboğan ve ark., 2008; Wolfe ve ark., 2015)]}

\kaynak{Arıboğan, D.Ü., Ortaylı, İ. ve Yavuz, H. (2008). \textit{Cumhuriyetimize Dair}. Aşiyan Yayınevi.}

\kaynak{Wolfe, J. M., Kluender, K. R., Dennis, L. M., Bartoshuk, L. M., Herz, R. S., Lederman, S. J. and Merfeld, D. M. (2015). \textit{Sensation \& perception }(4th ed.). Sinauer Associates.}

\kaynak{\textbf{Kitap Bölümü}}

\kaynak{\notrenk{FF0000}{Yazarın soyadı, Yazarın adının baş harfi. ve Yazarın soyadı, Yazarın adının baş harfi. (Yıl). Kitap bölümünün adı. }\notrenk{FF0000}{\textit{Kitabın adı}} \notrenk{FF0000}{(Baskı sayısı, Sayfa aralığı) içinde. Yayınevi.}}

\kaynak{[Metin içinde atıf: (Hovardaoğlu, 2007; Luck, 2014)]}

\kaynak{Hovardaoğlu, S. (2007). Psikolojik ölçmenin temelleri. Davranış bilimleri için araştırma teknikleri (2. Baskı, ss. 87-125) içinde. Hatipoğlu Yayınevi.}

\kaynak{Luck, S. J. (2014). A Broad Overview of the Event-Related Potential Technique. In An introduction to event related potential technique (2nd ed., pp. 1-34). The MIT Press.}

\kaynak{\textbf{Kitap: Farklı baskılar }}

\kaynak{\notrenk{FF0000}{Yazarın soyadı, Yazarın adının baş harfi. (Yıl). }\notrenk{FF0000}{\textit{Kitabın adı }}\notrenk{FF0000}{(Baskı Sayısı). Yayınevi.}}

\kaynak{[Metin içinde atıf: (Jackson, 2019; Tarhan, 2017; Wolfe ve ark., 2015)]}

\kaynak{Jackson, L. M. (2019). \textit{The psychology of prejudice: From attitudes to social action} (2nd ed.). American Psychological Association. https://doi.org/10.1037/0000168-000}

\kaynak{Tarhan, N. (2017). \textit{Toplum psikolojisi : sosyal şizofreniden toplumsal empatiye }(10. Baskı). Timaş Yayınları.}

\kaynak{Wolfe, J. M., Kluender, K. R., Dennis, L. M., Bartoshuk, L. M., Herz, R. S., Lederman, S. J. ve Merfeld, D. M. (2015). \textit{Sensation \& perception }(4th ed.). Sinauer Associates.}

\kaynak{\textbf{Kitap: Yazarı olmayan }}

\kaynak{[Metin içinde atıf: (Türkçe bilim terimleri sözlüğü, 2011)]}

\kaynak{\textit{Türkçe bilim terimleri sözlüğü : Sosyal bilimler}. (2011). Akademi Kitabevi.}

\kaynak{\textbf{Kitap: Editörlü }}

\kaynak{\notrenk{FF0000}{Editörün soyadı, editörün adının baş harfi. (Ed.). (Yıl). }\notrenk{FF0000}{\textit{Kitabın adı}} \notrenk{FF0000}{(Baskı sayısı). Yayınevi.}}

\kaynak{[Metin içinde atıf: (Hygum ve Pedersen, 2010; Kesharwani, 2020; Tanrıdağ, 2016)]}

\kaynak{Hygum, E., ve Pedersen, P. M. (Eds.). (2010). \textit{Early childhood education: Values and practices in Denmark}. Hans Reitzel Forlag.}

\kaynak{Kesharwani, P. (Ed.). (2020). \textit{Nanotechnology based approaches for tuberculosis treatment}. Academic Press}

\kaynak{Tanrıdağ, O. (Ed.). (2016). \textit{Davranış nörolojisi : Beyin-davranış ilişkilerinin organizasyon prensipleri, sendromları ve hastalıkları}. Nobel Tıp Kitapları}

\kaynak{\textbf{Editörlü Kitapta Bölüm}}

\kaynak{\notrenk{FF0000}{Yazarın soyadı, Yazarın adının baş harfi. (Yıl). Kitap bölümünün adı. Editörün adının baş harfi. Editörün soyadı (Ed.), }\notrenk{FF0000}{\textit{Kitabın adı}} \notrenk{FF0000}{(Baskı sayısı, Sayfa aralığı) içinde. Yayınevi.}}

\kaynak{[Metin içinde atıf: (Özkan Ceylan ve Bekçi, 2012; Posner, 1995)]}

\kaynak{Özkan Ceylan, A. ve Bekçi, B. (2012). Algı. N. Güngör Ergan, B. Şahin Kütük ve R. Coştur (Ed), \textit{Davranış bilimleri} (ss. 33-53) içinde. Siyasal Kitabevi}

\kaynak{Posner, M. I. (1995). Attention in cognitive neuroscience: An overview. In M. S. Gazzaniga (Ed.), \textit{The cognitive neurosciences} (pp. 615-624). The MIT Press}

\kaynak{\textbf{Kitap: Yazarı kurum olan}}

\kaynak{[Metin içinde atıf: (Üsküdar Üniversitesi, 2018)]}

\kaynak{Üsküdar Üniversitesi Yayınları. (2018). \textit{Yükseköğretim ve demokratik kültür : Vatandaşlık, insan hakları ve sivil sorumluluk, 30 Şubat 2018 }. Üsküdar Üniversitesi}

\kaynak{\textbf{Kitap: Çeviri }}

\kaynak{\notrenk{FF0000}{Orijinal kitabın yazarının soyadı, Adının baş harfi. (Yıl). }\notrenk{FF0000}{\textit{Kitabın adı}} \notrenk{FF0000}{(Baskı sayısı). (Çevirmenin adının baş harfi. Çevirmenin soyadı, Çev.). Yayınevi. (Orijinal eserin yayın tarihi )}}

\kaynak{[Metin içinde atıf: (Arasteh, 2000)]}

\kaynak{Arasteh, A. R. (2000). \textit{Aşkta ve yaratıcılıkta yeniden doğuş : Mevlana Celaleddin Rumi’nin kişilik çözümlemesi}. (Çev. B. Demirkol, ve İ. Özdemir,). Kitabiyat.}

\kaynak{\notrenk{0070C0}{\textbf{DERGİ MAKALELERİ}}}

\kaynak{\textbf{Basılı dergi makalesi }}

\kaynak{\notrenk{FF0000}{Yazarın soyadı, Yazarın Adının Baş Harfi. (Yıl). Makalenin başlığı. }\notrenk{FF0000}{\textit{Dergi adı}}\notrenk{FF0000}{\textit{ (İtalik)}}\notrenk{FF0000}{, Cilt(sayı), Sayfa aralığı. http://doi.org/xx.xxxxxxxxxx}}

\kaynak{[Metin içinde atıf: (Grady ve ark., 2019; Metin ve ark., 2018; Tarhan, 2001)}

\kaynak{Alper, E., ve Nevzat, T. (2020). Inflammation Biomarkers in Psychiatry. \textit{Current Psychiatry Research and Reviews}, \textit{16}(2), 78-85. https://doi.org/10.2174/2666082216999200625115701}

\kaynak{Grady, J. S., Her, M., Moreno, G., Perez, C., \& Yelinek, J. (2019). Emotions in storybooks: A comparison of storybooks that represent ethnic and racial groups in the United States. \textit{Psychology of Popular Media Culture}, 8(3), 207–217. https://doi.org/10.1037/ppm0000185}

\kaynak{Kalnay, E., Kanamitsu, M., Kistler, R., Collins, W., Deaven, D., Gandin, L., Iredell, M., Saha, S., White, G., Woollen, J., Zhu, Y., Chelliah, M., Ebisuzaki, W., Higgins, W., Janowiak, J., Mo, K. C., Ropelewski, C., Wang, J., Leetmaa, A., \ldots{} Joseph, D. (1996). The NCEP/NCAR 40-year reanalysis project. \textit{Bulletin of the American Meteorological Society}, 77(3), 437–471. http://doi.org/.....}

\kaynak{Metin, S. Z., Erguzel, T. T., Ertan, G., Salcini, C., Kocarslan, B., Cebi, M., Metin, B., Tanridag, O., \& Tarhan, N. (2018). The Use of Quantitative EEG for Differentiating Frontotemporal Dementia From Late-Onset Bipolar Disorder. \textit{Clinical EEG and Neuroscience}, 49(3), 171-176. https://doi.org/10.1177/1550059417750914}

\kaynak{Tarhan, N. (2001). Loneliness and Social Dissatisfaction in Turkish Adolescents. \textit{The Journal of Psychology}, \textit{135}(1), 113-123. https://doi.org/10.1080/00223980109603684}

\kaynak{\textbf{Elektronik dergi makalesi: Ön baskı }}

\kaynak{[Metin içinde atıf: (Hampton ve ark., 2017; Hetland ve ark., 20180)]}

\kaynak{Hetland, B., McAndrew, N., Perazzo, J., \& Hickman, R. (2018). A qualitative study of factors that influence active family involvement with patient care in the ICU: Survey of critical care nurse\textit{s}. \textit{PubMed Central}. https://www.ncbi.nlm.nih.gov/pmc/articles/PMC5736422/?report=classic}

\kaynak{Hampton, S., Rabagliati, H., Sorace, A., \& Fletcher-Watson, S. (2017). Autism and bilingualism: A qualitative interview study of parents’ perspectives and experience\textit{s}. \textit{PsyArXiv}. https://doi.org/10.31234/osf.io/76xfs}

\kaynak{\notrenk{0070C0}{\textbf{SÖZLÜK/ANSİKLOPEDİ}}}

\kaynak{\textbf{Sözlük ve Ansiklopedi}}

\kaynak{\notrenk{FF0000}{Yazarın soyadı, Adının baş harfi. (Yıl). }\notrenk{FF0000}{\textit{Kitabın adı}} \notrenk{FF0000}{(Baskı sayısı). Yayınevi.}}

\kaynak{[Metin içinde atıf: (Bilgin, 2003; Keyormarsi ve ark., 2020; VandenBos, 2013)]}

\kaynak{Bilgin, N. (2003). Sosyal psikoloji sözlüğü: Kavramlar, yaklaşımlar. Bağlam Yayıncılık.}

\kaynak{Keyormarsi, K., O’Leary, N., ve Pardee, A. B. (2020). \textit{Cell division}. McGraw-Hill Education. https://doi.org/10.1036/1097-8542.116300}

\kaynak{VandenBos, G. R. (Ed.). (2013). \textit{APA dictionary of clinical psychology}. American Psychological Association.}

\kaynak{\notrenk{0070C0}{\textbf{WEB KAYNAKLARI}}}

\kaynak{\textbf{Web Sayfası  }}

\kaynak{\notrenk{FF0000}{Yazar Soyadı, Adının baş harfi veya Grup adı. (Alıntı tarihi).Çalışmanın başlığı. İnternet sitesinin adı. URL adresinden .... tarihinde alınmıştır.}}

\kaynak{[Metin içinde atıf: (Türk Psikologlar Derneği, 2019; World Health Organization [WHO], 2020)]}

\kaynak{Türk Psikologlar Derneği (2019, 26 Kasım). Mesleki mevzuat. https://www.psikolog.org.tr/tr/kurumsal/meslekimevzuat-x654/ adresinden 2 Eylül 2020 tarihinde alınmıştır.}

\kaynak{World Health Organization (2020, 2 Eylül). Coronavirus. https://www.who.int/healthtopics/coronavirus\#tab=tab\_1}

\kaynak{\notrenk{0070C0}{\textbf{TEZLER}}}

\kaynak{\notrenk{FF0000}{Yazarın soyadı, Yazarın adının baş harfi. (Yıl). }\notrenk{FF0000}{\textit{Tezin baş}}\notrenk{FF0000}{lığı [Yüksek lisans tezi / Doktora tezi]. Üniversitenin adı.}}

{\teksatir [Metin içinde atıf: (Çetinkaya, 2015; Harris 2014)]\par}

\kaynak{Çetinkaya, Ş. (2015). \textit{Stochastic mortality using non-life methods} [Doktora tezi] Doğuş Üniversitesi}

\kaynak{Harris, L. (2014). \textit{Instructional leadership perceptions and practices of elementary school leaders} [Doktora tezi]. University of Virginia.}

{\teksatir Kabir, J. M. (2016). \textit{Factors influencing customer satisfaction at a fast food hamburger chain: The relationship between customer satisfaction and customer loyalty} [Doktora Tezi] Wilmington University\par}

{\teksatir Miranda, C. (2019). \textit{Exploring the lived experiences of foster youth who obtained graduate level degrees: Self-efficacy, resilience, and the impact on identity development} [Doktora tezi] Pepperdine University\par}

{\teksatir \notrenk{0070C0}{\textbf{KONFERANS BILDIRILERI }}\par}

{\teksatir \textbf{Dergide Yayınlanan Konferans Bildirileri}\par}

\kaynak{\notrenk{FF0000}{Yazarın soyadı, Yazarın Adının Baş Harfi. (Yıl). Bildirinin başlığı. }\notrenk{FF0000}{\textit{Derginin Adı}}\notrenk{FF0000}{, Cilt(Süreli yayının sayısı), Sayfa aralığı. http://doi.org/xx.xxxxxxxxxx}}

\kaynak{[Metin içinde atıf: (Duckworth ve ark., 2019)]}

\kaynak{Duckworth, A. L., Quirk, A., Gallop, R., Hoyle, R. H., Kelly, D. R., ve Matthews, M. D. (2019). Cognitive and noncognitive predictors of success. \textit{Proceedings of the National Academy of Sciences}, USA, 116(47), 23499–23504. https://doi.org/10.1073/pnas.1910510116}

{\teksatir \textbf{Yayımlanmamış konferans bildirisi}\par}

\kaynak{\notrenk{FF0000}{Yazarın soyadı, Yazarın adının baş harfi. (Yıl, Gün, Ay). Sözlü/poster sunumun adı [Sözlü/poster Sunum]. Bilimsel Toplantının Adı, Toplantının gerçekleştiği şehir, Ülke.}}

\kaynak{Özkan A. ve Şenyüz L. (2004, Eylül, 7-11). Haloperidolün, farklı uyaranlar tarafından kontrol edilen, su elde etmeye yönelik edimsel davranış üzerindeki etkileri [Sözlü sunum]. 13. Ulusal Psikoloji Kongresi, İstanbul, Türkiye.}

\kaynak{Salman, F. ve Cangöz, B. (2016, Eylül, 5-7). Alzheimer tipi demansta iç kaynaklı ve dış kaynaklı dikkat: duygusal bağlam etkisi [Poster sunumu]. 19. Ulusal Psikoloji Kongresi, İzmir, Türkiye.}

{\teksatir \textbf{Basılmış/Elektronik konferans kitabında bildiri }\par}

\kaynak{\notrenk{FF0000}{Yazarın soyadı, Yazarın adının baş harfi. (Yıl, Gün, Ay). Bildirinin adı. Editörün adının baş harfi. Editörün soyadı (Ed.), }\notrenk{FF0000}{\textit{Sempozyum Kitabının adı}} \notrenk{FF0000}{(Baskı sayısı, Sayfa aralığı) içinde. Toplantının gerçekleştiği şehir, Ülke}}

\kaynak{Altürk, S., Avci, D., Tamer, Ö. ve Atalay, Y. (2017, 6-10, Eylül). The second− and third−order nonlinear optical properties and Electronic transitions of a nlo chromophore: a DFTstudy. B. Akkuş, Y. Öktem, G. Süsoy, Ö. Aytan, A. Kurt, G. Aydoğdu, O. A.Barut (Ed.), \textit{9 th International Physics Conference of the Balkan Physical Union} içinde. İstanbul, Türkiye}

{\teksatir \notrenk{0070C0}{\textbf{RAPORLAR }}\par}

{\teksatir \textbf{Resmi raporlar }\par}

{\teksatir \notrenk{FF0000}{Yazarın soyadı, Yazarın Adının Baş Harfi. (Yıl). }\notrenk{FF0000}{\textit{Raporun}}\notrenk{FF0000}{\textit{ }}\notrenk{FF0000}{\textit{adı}} \notrenk{FF0000}{(Yayın no.).Yayınevi/Erişim adresi}\par}

{\teksatir \notrenk{FF0000}{Resmi Yayını Basan Kurumun Adı. (Yıl). }\notrenk{FF0000}{\textit{Raporun}}\notrenk{FF0000}{\textit{ }}\notrenk{FF0000}{\textit{adı}} \notrenk{FF0000}{(Yayın no.).Yayınevi/Erişim adresi}\par}

{\teksatir [Metin içinde atıf: (Akbaytürk Çanak ve ark., 2014, s. 40; Türkiye İstatistik Kurumu [TÜİK], 2019, s. 67)]\par}

\kaynak{Akbaytürk Çanak, T., Çelik, S., Çetinkaya, İ., Çukadar, S., Güneş, G., Gürdal, G., Holt İ., Kaya E. ve Kaygusuz, A. (2014). \textit{2023’e doğru Türkiye’de üniversite kütüphaneleri: Mevcut durum, sorunlar, standartlar ve çözüm önerileri}. Yükseköğretim Kurulu. Erişim adresi https://yok.gov.tr/web/ukr}

\kaynak{Türkiye İstatistik Kurumu. (2019). \textit{İstatistiklerle çocuk} (Yayın no.4581). https://biruni.tuik.gov.tr/yayin/views/visitorPages/index.zul}

{\teksatir \notrenk{0070C0}{\textbf{PATENT }}\par}

{\teksatir \notrenk{FF0000}{Yazarın soyadı, Yazarın Adının Baş Harfi. (Yıl). }\notrenk{FF0000}{\textit{Patentin}}\notrenk{FF0000}{\textit{ }}\notrenk{FF0000}{\textit{adı}}\notrenk{FF0000}{. Patent Ülkesi. Patent No}\par}

{\teksatir [Metin içinde atıf: (Apelsmeier, 2013)]\par}

{\teksatir Apelsmeier, A. (2013).\textit{ Motor vehicle having an air conditioner compressor as a starter of the internal combustion engine. }U.S. Patent No. US9731704B2\textit{ }\par}

\begin{kaynaklistesi}
  \item Abrahart, R. J. ve See, L. (1998, 17–19, Eylül). Neural Network vs. ARMA Modelling: Constructing Benchmark Case Studies of River Flow Prediction. In GeoComputation ’98. Proceedings of the Third International Conference on GeoComputation, University of Bristol, United Kingdom

  \item Abrahart, R. J. ve See, L. (2000). Comparing neural network and autoregressive moving average techniques for the provision of continuous river flow forecasts in two contrasting catchments, \textit{Hydrological Processes,} 14, 2157–2172.

  \item Spielberg, S. (Producer) ve Spielberg, S. (Director). (1993). \textit{Schindler’s list }[Motion picture]. California : Universal Picture.

\end{kaynaklistesi}


\ekler
\textbf{EK A. }Haritalar

\textbf{EK B. }Tablolar

\clearpage

\ek{Haritalar}

\begin{figure}[H]
  \centering
  \includegraphics[width=11.11cm]{sekiller/image1.png}
  \caption{İklim haritaları: (a)Avustralya. (b) Afrika. (c)Antartika. (d). \ldots{}}
\end{figure}

\clearpage

\begin{table}[H]
  \caption{Ekler bölümünde tablo örneği.}
  \centering
\begin{tabular}{C{3.05cm}C{2.10cm}C{2.16cm}C{2.52cm}}
\hline\hline
\Ns Kolon A & Kolon B & Kolon C & Kolon D \\
\hline
\Ns Satır A & Satır A & Satır A & Satır A \\
\Ns Satır B & Satır B & Satır B & Satır B \\
\Ns Satır C & Satır C & Satır C & Satır C \\
\hline
\end{tabular}
\end{table}

\clearpage


\begin{ozgecmis}
% [Altyaz] sp=before:120
Ad-Soyad\quad{}\quad{}  :    Xxxxxxx XXXXXXXX

% [Normal] sp=line:240,lineRule:auto
\textbf{ÖĞRENİM DURUMU:}

\begin{itemize}
  \item \textbf{Lisans}\quad{}\textbf{            :} Mezuniyet yılı, Üniversite, Fakülte, Bölüm

  \item \textbf{Yükseklisans       :} Mezuniyet yılı, Üniversite, Anabilim Dalı, Program

\end{itemize}
% [Normal] sp=line:240,lineRule:auto
\textbf{MESLEKİ DENEYİM VE ÖDÜLLER:}

\begin{itemize}
  \item 2016-2017 yılları arasında Milli Eğitim Bakanlığı’nda öğretmen olarak çalıştı.

  \item 2020 yılında Sakarya Üniversitesi Dönemsel Başarı Ödülü’nü kazandı.

\end{itemize}
% [Normal] sp=line:240,lineRule:auto
\textbf{TEZDEN TÜRETİLEN YAYINLAR/SUNUMLAR}

\begin{itemize}
  \item Ganapuram S., Hamidov A., Demirel, M. C., Bozkurt E., Kındap U. ve Newton A. (2007, 22-24, Mart). Erasmus Mundus Scholar's Perspective On Water And Coastal Management Education In Europe. \textit{International Congress - River Basin Management}, Antalya, Turkey. (\notrenk{FF0000}{Sunum örneği)}

  \item Satoğlu, Ş.I., Durmuşoğlu, M. B. ve Ertay, T. A. 2010. A Mathematical Model And A Heuristic Approach For Design Of The Hybrid Manufacturing Systems To Facilitate One-Piece Flow, \textit{International Journal of Production Research}, 48(17), 5195-5220\notrenk{FF0000}{.  (Makale örneği})

  \item Chen, Z. 2013. Intelligent Digital Teaching And Learning All-In-One Machine, Has Projection Mechanism Whose Front End Is Connected With Supporting Arm, And Base Shell Provided With Panoramic Camera That Is Connected With Projector. Patent numarası: CN203102627-U  (\notrenk{FF0000}{Patent örneği})

\end{itemize}
% [Normal] sp=line:240,lineRule:auto
\textbf{DİĞER ESERLER:}
\end{ozgecmis}

\end{document}
